\section{AMC 12 12A 2025}
        \begin{enumerate}
        	\item (\textit{AMC 12 12A 2025}) A souvenir shop sells solid metal sculptures of Achilles in two sizes. One is $6$ cm tall and one is $8$ cm tall. Proportions between the two sculptures are preserved. The smaller sculpture has a mass of $486$ grams. What is the mass, in grams, of the larger sculpture?


 $\textbf{(A)}\   648  \qquad                 \textbf{(B)}\   864  \qquad                 \textbf{(C)}\   972  \qquad                 \textbf{(D)}\   1152  \qquad                 \textbf{(E)}\   1458$

 \textbf{Answer: }E



\textbf{Solution 1}

At $2{:}30$, Andy is $8$ miles ahead. For every hour that they both travel, Betsy gains $4$ miles on Andy. Therefore, it will take her $2$ more hours to be the same distance from the starting point as Andy, which occurs at $\boxed{\textbf{(E) } 4{:}30}$.

 $8h = 12(h-1)$ $\Rightarrow 8h = 12h - 12$ $\Rightarrow 4h = 12$ $\Rightarrow h = 3$

 Since Andy started at $1{:}30$, the catch-up time is $4{:}30$.  Answer: $\textbf{(E)}$.

 Alternatively, from Betsy's perspective:  $8(h+1) = 12h$ $\Rightarrow 8h + 8 = 12h$ $\Rightarrow h = 2$  Same result: $\textbf{(E) } 4{:}30$.

 



\textbf{Solution 2 (Simplest Version)}

Andy goes $8$ mph, Betsy goes $12$ mph.  

 We know that Betsy starts 1 hour after Andy, so she will be 8 miles behind.

 $12x - 8 = 8x$
$4x = 8$ $\,\rightarrow\,$ $x = 2$

 2 hours after Betsy's start time is $\boxed{\text{4:30}}$

 -XaberTooth

 -CMM

 



\textbf{Solution 3}

Try the answer choices:

 At $3{:}30$: Andy $=16$, Betsy $=12$  At $3{:}45$: Andy $=18$, Betsy $=15$  At $4{:}00$: Andy $=20$, Betsy $=18$  At $4{:}15$: Andy $=22$, Betsy $=21$  At $4{:}30$: Andy $=24$, Betsy $=24$You can visualize this using a table:

 \begin{array}{c|c|c}\textbf{Time} & \textbf{Andy - 8 mph} & \textbf{Betsy - 12 mph} \\ \hline3{:}30 & 16 & 12 \\3{:}45 & 18 & 15 \\4{:}00 & 20 & 18 \\4{:}15 & 22 & 21 \\4{:}30 & 24 & 24\end{array}

 Thus the correct answer is $\boxed{\textbf{(E) } 4{:}30}$.

 



\textbf{Solution 4}

Since $\text{lcm}(8,12)=24$, the time for Andy to reach 24 miles is  $24/8 = 3$ hours.

 $1{:}30 + 3{:}00 = \boxed{\textbf{(E) } 4{:}30}$.

 Note: Alternatively we can find that the time for Betsy to reach $24$ miles is $\frac{24}{12} = 2$ hours. So we have

 
 \[2{:}30 + 2{:}00 = \boxed{\textbf{(E) } 4{:}30}.\]




\textbf{Solution 5}

Andy goes $8x$ miles, Betsy goes $12(x-1)$ miles.  Set equal:

 $8x = 12(x-1)$ $\Rightarrow x = 3$

 Thus the answer is $\boxed{\textbf{(E) } 4{:}30}$.

 



\textbf{Solution 6}

Using constant rates:

 
 \[\begin{array}{c|c|c} \textbf{Time} & \textbf{Andy} - 8\text{ mph} & \textbf{Betsy} - 12\text{ mph} \\ \hline \hline 1{:}30 - \text{Starting} & 0 \text{ miles} & 0 \text{ miles} \\ \hline 2{:}30 & 8 \text{ miles} & 0 \text{ miles} \\ \hline 3{:}30 & 16 \text{ miles} & 12 \text{ miles} \\ \hline  4{:}30 & 24 \text{ miles} & 24 \text{ miles} \end{array}\]

 Thus the answer is $\boxed{\textbf{(E) } 4{:}30}$.

 



\textbf{Solution 7 (Calculus Overkill)}

Let  $D_A(t) = 8t$,  $D_B(t) = 12(t-1)$ for $t \ge 1$.

 Difference:  $F(t) = 12(t-1) - 8t = 4t - 12$.

 Set $F(t)=0$:  $4t - 12 = 0$ $t = 3$ hours.

 $1{:}30 + 3 = 4{:}30$.

 -apex304

 



\textbf{Solution 8 (Easy to understand, fast)(Detailed)}

Let \( t \) denote time in hours measured from 1:30 PM:

 \( t := \text{hours after } 1:30 \text{ PM} \).

 Thus \( t = 0 \) at 1:30 PM, \( t = 1 \) at 2:30 PM, and the physically relevant domain for both riders having left is \( t \geq 1 \).

 Velocity/position (vectorized).Choose orthonormal basis \( \{ i, j \} \) with \( i \) east and \( j \) north. Define piecewise position vectors \( r_A(t), r_B(t) \in \mathbb{R}^2 \) (for all \( t \geq 0 \)):

 \( r_A(t) = 8t \, j = (0, 8t) \),
 \[r_B(t) = \begin{cases} (0, 0), & 0 \leq t < 1, \\ 12(t - 1) \, i, & t \geq 1, \end{cases} = (12 \max\{t - 1, 0\}, 0).\]

 Define scalar distance functions (Euclidean norm / 2-norm):

 \( D_A(t) := \| r_A(t) \|_2 = \sqrt{0^2 + (8t)^2} = 8t \),\( D_B(t) := \| r_B(t) \|_2 = 12 \max\{t - 1, 0\} \).

 We seek \( t \) with \( D_A(t) = D_B(t) \) and \( t \geq 1 \). For \( t \geq 1 \) the equality reduces to the linear equation

 \( 8t = 12(t - 1) \).

 Bring like terms together:

 \( 8t = 12t - 12 \implies 12 = 12t - 8t \).

 Compute the difference \( 12t - 8t \) digit-by-digit: \( 12t - 8t = (12 - 8)t = 4t \). So

 \( 12 = 4t \).

 Divide \( 12 \) by \( 4 \) carefully: \( 12 \div 4 = 3 \). Hence

 \( t = 3 \).

 \begin{verbatim}
\href{https://youtu.be/LN5ofIcs1kY}{https://youtu.be/LN5ofIcs1kY} 
\end{verbatim}

\begin{verbatim}
\href{https://youtu.be/c-UDo53KwFU}{https://youtu.be/c-UDo53KwFU}
\end{verbatim}

\textbf{Easy Solution} 


	\item (\textit{AMC 12 12A 2025}) Bees occupy the $8$ corners of a $1 \times 2 \times 3$ rectangular box. The bees can fly inside the box, but cannot fly farther than $\frac{1}{2}$ unit from their respective corners. What is the volume, in cubic units, of the region inside the box that cannot be reached by any bee?


 $\textbf{(A)}\   6 - \frac{\pi}{3}  \qquad                 \textbf{(B)}\   6 - \frac{\pi}{4}  \qquad                 \textbf{(C)}\   6 - \frac{\pi}{6}  \qquad                 \textbf{(D)}\   6 - \frac{\pi}{8}  \qquad                 \textbf{(E)}\   6 - \frac{\pi}{12}$

 \textbf{Answer: }B



\textbf{Solution 1}

We are given that $0.2(10) = 2$ pounds of cashews in the first box.

 Denote the pounds of nuts in the second nut mix as $x.$ The total peanuts in the total mix will be:

 
 \[\frac{5 + 0.2x}{10 + x} = 0.4\]

 Then solve for $x.$

 \[5 + 0.2x = 0.4(10 + x)\]
 \[5 + 0.2x = 4 + 0.4x\]
 \[5 - 4 = 0.4x - 0.2x\]
 \[0.2x = 1\]
 \[x = 5\]

 Thus, we have $5$ pounds of the second mix. 

 
 \[0.2(10) + 0.4(5) = 2 + 2 = \boxed{\text{(B) }4}\]




\textbf{Solution 2}

Let the number of pounds of nuts in the second nut mixture be $x$. Therefore, we get the equation $0.5 \cdot 10 + 0.2 \cdot x = 0.4(x+10)$. Solving it, we then get $x=5$. Therefore, the amount of cashews in the two bags is $(0.2 \cdot 10) + (0.4 \cdot 5) = 4$, so our answer choice is $\boxed{\textbf{(B)} 4}$.

 



\textbf{Solution 3}

The percent of peanuts in the first mix is $10\%$ away from the total percentage of peanuts, and the percent of peanuts in the second mix is $20\%$ away from the total percentage. This means the first mix has twice as many nuts as the second mix, so the second mix has $5$ pounds. $0.20 \cdot 10 + 0.40 \cdot 5 = 4$ pounds of cashews. So our answer is, $\boxed{\textbf{(B)}4}$

 The weight of the second box of nut mix can be calculated by: ($0.2x + 5) / (10 + x) = 0.4$, where the numerator is equal to the weight of the peanuts, the denominator is equal to the total weight of the boxes, and $x$ is equal to the weight of the second box.

 Multiplying both sides by $10 + x$, you get $0.2x + 5 = 4 + 0.4x$

 Isolating the variable: $0.2x = 1$$\Rightarrow x = 5$

 Calculating the weight of the cashews: $5(0.4) + 10(0.2)$$\Rightarrow 2 + 2$$\Rightarrow 4$

 The answer is $\boxed{\textbf{(B) } 4}$.

 



\textbf{Solution 4}

Note that we can set the information given in the problem into a table shown below:

 
 \[\renewcommand{\arraystretch}{1.5} \begin{centering} \begin{array}{| c | c | c |} \hline \text{Peanuts} & \text{Cashews} & \text{Almonds}\\ \hline 5 & 2 & 3\\ \hline \frac{2}{10}x & \frac{4}{10}x & \frac{4}{10}x\\ \hline \end{array} \end{centering}\]

 We are given that the new nut mix will contain $40\%$ peanuts. Hence, $5 + \frac{2}{10}x$ is $40\%$ of the total mix which is $10 + x$. Solving the equation $5 + \frac{2}{10}x = \frac{2}{5} \cdot (10 + x)$ yields $x=5.$ Therefore, the number of cashews in the new mix is equal to $2 + \frac{2}{5} \cdot 5 = \boxed{\textbf{(B)}  4}$.

 



\textbf{Solution 5}

Write composition vectors as column vectors (mass fractions multiplied by mass):

 Initial mass vector (peanuts, cashews, almonds) for 10 lb:
 \[\mathbf{v}_0 = 10 \begin{pmatrix} 0.50 \\ 0.20 \\ 0.30 \end{pmatrix} = \begin{pmatrix} 5 \\ 2 \\ 3 \end{pmatrix}.\]Added mix per unit mass:
 \[\mathbf{u} = \begin{pmatrix} 0.20 \\ 0.40 \\ 0.40 \end{pmatrix}.\]After adding \( x \) lb, total mass vector is
 \[\mathbf{v}(x) = \mathbf{v}_0 + x \mathbf{u} = \begin{pmatrix} 5 \\ 2 \\ 3 \end{pmatrix} + x \begin{pmatrix} 0.2 \\ 0.4 \\ 0.4 \end{pmatrix}.\]Peanut fraction requirement:
 \[\frac{\mathbf{e}_1^\top \mathbf{v}(x)}{\mathbf{1}^\top \mathbf{v}(x)} = 0.4,\]where \( \mathbf{e}_1 = \begin{pmatrix} 1, 0, 0 \end{pmatrix}^\top \) and \( \mathbf{1} = \begin{pmatrix} 1, 1, 1 \end{pmatrix}^\top \). This is exactly the same scalar equation,
 \[\frac{5 + 0.2x}{10 + x} = 0.4,\]leading to \( x = 5 \) and hence
 \[\mathbf{v}(5) = \begin{pmatrix} 5 \\ 2 \\ 3 \end{pmatrix} + 5 \begin{pmatrix} 0.2 \\ 0.4 \\ 0.4 \end{pmatrix} = \begin{pmatrix} 6 \\ 4 \\ 5 \end{pmatrix}.\]Thus cashews \( = 4 \) lb.

 



\textbf{Solution 6: Very easy algebra and straightforward!}

Set up the ratios for the initial and the additional mixtures. The ratio of Peanuts:Cashews:Almonds (denoted as P:C:A) would initially be $5x:2x:3x$ for some factor $x$. Then for the second mixture it would be $2x:4x:4x$. If we add these we get 7x:6x:7x. Now we know that the percentage of peanuts is 40\%. But we also know that the number of pounds of peanuts and almonds is the same. $40\%+40\%=80\%$ and so the number of cashews makes up $20\%$ of the mixture. If you divide the factor $x$ by the whole ratio you get $7:6:7$ . These values sum to $20$ and $20\%$ of $20$ is $4$.

 \textbf{Video Solutions} 


	\item (\textit{AMC 12 12A 2025}) There are $2026$ equally spaced points that all lie on a circle. A line segment is drawn between two points if there is exactly one point between them along the circle. Into how many regions will this divide the interior of the circle? The figure below shows an example with $10$ equally spaced points, in which case there are $21$ interior regions.


 
                    \begin{figure}[H]
                        \centering
                        \includegraphics[width=0.5\textwidth]{Problem/AMC_12_12A_2025_3_Problem_Figure_1.png}
                    \end{figure}
                    
$\textbf{(A)}\   4052  \qquad                 \textbf{(B)}\   4053  \qquad                 \textbf{(C)}\   4054  \qquad                 \textbf{(D)}\   6079  \qquad                 \textbf{(E)}\   6081$

 \textbf{Answer: }A



\textbf{Solution 1}

When Ash joins a team, the team's average is pulled towards his age. Let $A$ be Ash's age and $N$ be the number of people on the student team. This means that there are $15-N$ people in the teacher team. Let us write an expression for the change in the average for each team.

 The students originally had an average of $12$, which became $14$ when Ash joined, so there was an increase of $2$. The term $A-12$ represents how much older Ash is compared to the average of the students'. If we divide this by $N+1$, which is the number of people on the student team when Ash joins, we get the average change per team member once Ash is added. Therefore, 
 \[\frac{A-12}{N+1} = 2.\]

 Similarly, for teachers, the average was originally $55$, which decreased by $3$ to become $52$ when Ash joined. Intuitively, $55-A$ represents how much younger Ash is than the average age of the teachers. Dividing this by the expression $(15-N)+1$, which is the new total number of people on the teacher team, represents the average change per team member once Ash joins. We can write the equation

 
 \[\frac{55-A}{16-N} = 3.\]

 To solve the system, multiply equation (1) by $N+1$, and similarly multiply equation (2) by $16-N$. Then add the equations together, canceling $A$, leaving equation $43=50-N$. From this we get $N=7$ and $A= \boxed{28}.$




\textbf{Solution 2}

As shown above, we get the equation 
 \[\frac{55-A}{16-N} = 3.\] Rearranging we get 
 \[\frac{55-A}{3} = 16-N.\] Therefore, since N must be an integer, $55-A$ must be divisible by 3. Plugging in values gets us $A= \boxed{28}.$




\textbf{Solution 3}

Another way is to say that there are $S$ students and $T$ teachers and Ash's age is $A$. We can create the equations\begin{align}&(1) S+T=15 \\&(2) \frac{12S+A}{S+1}=14 \\ &(3) \frac{55T+A}{T+1}=52\end{align}Equation $(2)$ simplifies to $A=2S+14$ and equation $(3)$ simplifies to $A=-3T+52$. Multiply equation $(2)$ by $3$ to get equation \begin{align} (4) 3A=6S+42 \end{align} and equation $(3)$ by $-2$ to get \begin{align}(5) -2A=6T-104 \end{align} Add (4) and (5) to get $A=6(S+T)-62$. After substituting equation (1) and simplifying, you get $A=\boxed{28}$ or answer choice $\boxed{A}$.

 



\textbf{Solution 4 (Cheese)}

Because the change in the kids group is less, you can assume that either there are more kids or the number is closer to the kids. If you try out 28 (Smallest number), with 7 kids, you notice the change is greater than the change in the mean of the kids ages in the problem, meaning that there must be 8 kids. Now know that there are 8 kids, to change the age by 2 the age of Ash must be 16 more than the average (because it changes by one for every 8, due to there being 8 kids), meaning Ash's age is 12 + 16, or $\boxed{28}$ (which is answer choice $\boxed{A}$).~ToxicWaste

 


 


	\item (\textit{AMC 12 12A 2025}) For how many integers $n$ is

 \[\frac{n^{2}-1}{n+2}\]
also an integer?


 $\textbf{(A)}\   2  \qquad                 \textbf{(B)}\   3  \qquad                 \textbf{(C)}\   4  \qquad                 \textbf{(D)}\   6  \qquad                 \textbf{(E)}\   \text{infinitely many}$

 \textbf{Answer: }B



\textbf{Solution 1}

We first number all the statements:

 1) At least one of these statements is true.2) At least two of these statements are true.3) At least two of these statements are false.4) At least one of these statements is false.

 We can immediately see that statement 4 must be true, as it would contradict itself if it were false. Similarly, statement 1 must be true, as all the other statements must be false if it were false, which is contradictory because statement 4 is true. Since both 1 and 4 are true, statement 2 has to be true. Therefore, statement 3 is the only false statement, making the answer $\boxed{\text{(B) }1}$. -Rainjs

 



\textbf{Solution 2}

Statements $I,II,$ and $IV$ are true, while statement $III$ is false. Hence, there are $3$ true statements and $\boxed{\text{(B) }1}$ false statement. This result can be checked by examining the statements individually again.

 Statements $I$ and $II$ will be true because there are $3\ge2$ true statements. Statement $IV$ is also true because there is $1\ge1$ false statement. Finally, statement $III$ is false because there are $1\ngeq2$ false statements.

 Note: This solution does not give a path to the solution, but rather only how to check the answer after the answer has been identified.

 



\textbf{Solution 3}

Let's say there are $T$ true statements. We know that $T$ can be any integer from $0$ to $4$. Numbering the $4$ given statements in order as $S1$ to $S4,$ we can represent $S1$ as $T \geq 1$, $S2$ as $T \geq 2$, and since there are $(4-T)$ false statements, $S3$ as $4-T \geq 2 \iff T \leq 2$, and $S4$ as $4-T \geq 1 \iff T \leq 3$.

 Accordingly, we now consider the possible cases for the value of $T$. If $T=0$, then the true statements are $S3$ and $S4$, i.e. there are $2$ true statements in total. But that means $T = 2 \neq 0$, giving a contradiction.

 Similarly, if $T=1$, then the true statements are $S1$, $S3$, and $S4$, i.e. there are $3$ true statements in total. This means $T = 3 \neq 1$, again giving a contradiction.

 If $T=2$, then all $4$ statements are true, yielding the contradiction $T = 4 \neq 2$.

 Finally, however, if $T=3$, then the true statements are $S1$, $S2$, and $S4$. Hence there are again $3$ true statements in total, but now as $T=3$ in the first place, this is actually consistent and not a contradiction.

 For completeness, we also observe that if $T=4$, then the true statements are only $S1$ and $S2$, i.e. $T = 2 \neq 4$, and so this case again leads to a contradiction.

 Thus we conclude that the only possible value is $T = 3$, which gives a consistent solution when $S1$, $S2$, and $S4$ are true and $S3$ is false (as noted above). It follows that the number of false statements is $\boxed{\textbf{(B) }1}$, namely precisely $S3$ (and indeed, $T = 3$ gives $4-T = 4-3 = 1$, as expected).

 



\textbf{Solution 4}

Suppose Statement $I$ is false, then none of the Statements are true, which contradicts the fact that a false Statement $III$ or $IV$ is telling the truth. Therefore, Statement $I$ is true and assume Statement $II$ is false. 

 Statement $II$ thus implies that only Statement $I$ was the truth, and the rest, false. But then, there are 3 false statements but then Statement $III$ and Statement $IV$ are telling the truth. So Statement $II$ is also true. 

 Now, if Statement $III$ is true, then both Statement $III$ and Statement $IV$ is false, contradicting the fact that it is true. Statement $III$ is hence false and Statement $IV$ tells the truth since Statement $III$ lied so indeed, there are at least one lie. There are a total of 3 truths and 1 lie, making the answer $\boxed{\text{(B) }1}$. ~hxve

 



\textbf{Solution 5 (Quick)}

I tried setting that the statements regarding the true ones (I and II) were true and regarding the false ones (III and IV) were false, and it clearly doesn't work.It doesn't work because statement III was true, not false as I set it to be originally.

 What if we make that statement (Statement III) true too? Well, then we have 3 T and 1 F statements, so statement IV is false, statement III is true, and statements I and II are true too.

 So there is $\boxed{\text{(B) }1}$ false statement.

 



\textbf{Solution 6 (Plug in Answer Choices)}

For this question, what I did was to simply test each of the answer choices. 

 Let's start with $\text{(A)}$ $0$ false statements:

 We can check the answer of $0$ false statements against each of the statements.

 Statement $\mathrm{I}$: "At least one of these statements is true."

 This statement is false, which already invalidates the answer of (A) 0 false statements, as we already have a false statement. 

 Let's move onto $\text{(B)}$ $1$ false statement:

 We can check once again this answer against each of the statements.

 Statement $\mathrm{I}$: "At least one of these statements is true."

 This statement is true, as when we have $1$ false statement, we have $3$ true statements, and therefore this is true.

 Statement $\mathrm{II}$: "At least two of these statements are true."

 Once again, this statement is true, as when we have $1$ false statement, we have $3$ true statements, and thus this is also true.

 Statement $\mathrm{III}$: "At least two of these statements are false."

 This statement is false, as we have one false statement, not $2$ and above.

 Statement $\mathrm{IV}$: "At least one of these statements is false."

 This statement is true, as we have one false statement.

 So, after checking the answer choice of $\boxed{\text{(B)} 1}$ false statement across the $4$ statements, we get three true statements and one false statement, making this statement the correct answer. I will not dive into the other $3$ answer choices, as we have already gotten the answer for this question, and if you were doing this way of solving on the test, you would most likely not test the other $3$ answers either for the sake of saving time. 

 



\textbf{Solution 7 (Step by Step, Simple)}

Start by going statement by statement and checking if they can be false. If statement 1 were false statement 4 would be true leading to a contradiction so statement 1 is true. If statement 2 were false 4 would be true and one would be true making statement 2 true. If statement 3 was false and we know statement 1 and statement 2 were true and 4 can be true and it works. Therefore, the correct answer is $\boxed{\text{(B) }1}$.

 --TFEA(would also appreciate if somebody could please latex my proof).--Iceboy1221(i wrote it using latex)

 



\textbf{Solution 8 (for Chinese Committee Paper)}

Due to reasons to prevent cheating, there was a variant for test takers in China, which involves slight iterations in the questions and answer choices. For this particular problem, the question asks for the number of $\mathbf{true}$ statements that Agnes wrote.

 Therefore, utilizing any of the methods in the solutions above, we find that there is one false statement, therefore there are three true statements.

 Therefore, the correct answer choice is $\mathbf{3}$.

 



\textbf{Solution 9}

Call the statements A, B, C, and D. We start by looking at statement C, since it forces some other things. If statement C is true, then statement D is necessarily also true. But there are two or more false statements, so that means both A and B would be false. But if A is false, there are no true statements, which is wrong, so statement C must be false.

 Since statement C is false, there can only be zero or one false statements. Since C is already false, this forces A, B, and D to all be true. This means that there is $\boxed{\text{(B) } 1}$ false statement.

 


	\item (\textit{AMC 12 12A 2025}) Let $a = \log_{2}27$, let $b = \log_{3}25$, and let $c = \log_{32}5$. What is $\frac{ab}{c}$?


 $\textbf{(A)}\   25  \qquad                 \textbf{(B)}\   27  \qquad                 \textbf{(C)}\   30  \qquad                 \textbf{(D)}\   32  \qquad                 \textbf{(E)}\   36$

 \textbf{Answer: }D



\textbf{Solution 1}

Let the side length of the largest square be $a,$ so it has area $a^2.$ Hence, the second-largest square has area $a^2k^2,$ the third-largest has $a^2k^4,$ and so on. 

 It follows that the total shaded area is 
 \[a^2-a^2k^2+a^2k^4-a^2k^6+...=a^2(1-k^2+k^4-k^6+...)=a^2\dfrac{1}{1+k^2}.\]

 The ratio of the area of the shaded region to that of the original square is then 
 \[\dfrac{a^2\frac{1}{1+k^2}}{a^2}=\dfrac{1}{1+k^2}=\dfrac{64}{100}\]
 \[\implies 64+64k^2=100\implies k^2=\dfrac{36}{64}\implies k=\boxed{\text{(D) }\dfrac{3}{4}}.\]

 \textbf{Remark} We can just let $a=1$ because the question deals with ratios, meaning that there wouldn't be a loss of generality if we let the side length equal some value, getting the same answer $\boxed{D}$. This was done in solution 3.

 



\textbf{Solution 2}

Let the side length of the first square be $a$ and the second square be $b$. The area of the original square is $a^2.$ The area of the outermost shaded region is $a^2 - b^2.$ Let $\frac{a}{b} = r.$ We have $b = \frac{a}{r}.$ So the outermost shaded region area becomes: 

 
 \[a^2 - b^2 = a^2 - \frac{a^2}{r^2}\]

 Now let the next square's side lengths be $c$ and $d.$ Similarly, $c = \frac{b}{r} = \frac{a}{r^2}$ and $d = \frac{c}{r} = \frac{a}{r^3}$and the area of the next shaded region becomes: 
 \[c^2 - d^2 = \frac{a^2}{r^4} - \frac{a^2}{r^6}\] 

 Notice the pattern of adding $4$ to the exponent. If this sequence continues infinitely, we ultimately get: 

 
 \[(a^2 + \frac{a^2}{r^4} + \frac{a^2}{r^8} + ...) - (\frac{a^2}{r^2} + \frac{a^2}{r^6} + ...)\] 

 Which can be simplified using the infinite geometric sequence formula. The problem also tells us that all of this equals $64\% a^2,$ or $\frac{16}{25}a^2:$

 \[\frac{a^2}{1 - \frac{1}{r^4}} - \frac{\frac{a^2}{r^2}}{1 - \frac{1}{r^4}} = \frac{16}{25}a^2\]

 \[\frac{1}{1 - \frac{1}{r^4}} - \frac{\frac{1}{r^2}}{1 - \frac{1}{r^4}} = \frac{16}{25}\]

 \[r^2 = \frac{16}{9} \implies r = \frac{4}{3}\] Since the problem tells us the ratio must be less than one, we take the reciprocal to finally get  
 \[k = \frac{1}{r} = \boxed{\text{(D) }\dfrac{3}{4}}.\]




\textbf{Solution 3 (Faster)}

Let the outside square be of side length $1$, and let $p = k^2$. Then the area of the square is $1-p+p^{2}-p^{3}+p^{4}-p^{5} \dots$ or $\left(1-p\right)\left(1+p^{2}+p^{4}+p^{6}+\dots\right)$. Then, using geometric series, we get $\left(1-p\right)\left(\frac{1}{1-p^{2}}\right) = \frac{1}{1+p}$. This is equal to $\frac{64}{100}$. Therefore, $1+p=\frac{25}{16}$, so $p = 9/16$ and $k = \frac{3}{4}$. Therefore, the answer is $\boxed{\text{(D)}\frac{3}{4}}.$




\textbf{Solution 4 (The fastest, no sequences/series)}

Since this is an infinite fixed pattern, we can actually take a huge shortcut. Suppose we take off the outer "layer" and then swap the colors. We have essentially created a smaller copy of the original diagram, specifically with a $k:1$ ratio of side lengths. Assume the original diagram had 100 units of area. Then, based on the givens, the original diagram had 64 units of shaded area and 36 units of unshaded area, meaning our new smaller diagram has 36 units of shaded area. Since the ratio of corresponding side lengths between our new diagram and the original diagram is $k$, the ratio of corresponding areas is $k^2$. Thus, we get $k^2 = \frac{36}{64}$, so $k = \boxed{\text{(D) }\dfrac{3}{4}}$.

 Note: the explanation seems complicated to explain, but the thought process while doing it is extremely fast, much faster than any of the methods involving a series.

 


 



\textbf{Solution 5 (based very closely on solution 4)}

Similarly to solution 4, we can imagine "peeling" off the outer layer to get an inner figure that is similar to the outer figure. Letting the large square have side 1, we know the outer shaded layer has area $1-k^2$. The shaded part of the inner figure (with the outer shaded part removed now must make up $\frac{36}{100} = \frac{9}{25}$ of the inner area of $k^2$. Together, the inner shaded part and the outer shaded part must add up to $\frac{64}{100}$, as they make up 64\% of the square. Thus, we get the equation$1-k^2 + \frac{9k}{25} = 16/25$- Solving this gives k = 3/4~Neber100

 



\textbf{Solution 6 (The longest way)}

Notes that: The sum $S$ can be written as:
 \begin{align*}     S &= 1^2 - k^2 + (k^2)^2 - (k^3)^2 + (k^4)^2 - (k^5)^2 + (k^6)^2 - \dots \\     &= k^0 - k^2 + k^4 - k^6 + k^8 - k^{10} + \dots \\     &= (k^0 + k^4 + k^8 + k^{12} + \dots) - (k^2 + k^6 + k^{10} + k^{14} + \dots) \\ \end{align*}
So we have: 
 \begin{align*}     S &= \lim_{n\rightarrow\infty}\left[1\times\dfrac{1-(k^4)^n}{1-k^4}-k^2\times\dfrac{1-(k^4)^n}{1-k^4}\right]\\     &= \dfrac{1}{1-k^4}-\dfrac{k^2}{1-k^4} \\     &= \dfrac{1-k^2}{1-k^4} \\     &= \dfrac{1-k^2}{(1+k^2)(1-k^2)} \\     &= \dfrac{1}{1+k^2} \end{align*}
We are given the condition that:
 \[\because \dfrac{1}{1+k^2}=64\%=\dfrac{64}{100}=\dfrac{16}{25}\]We can now solve for $k$:
 \begin{align*}   \therefore \dfrac{1}{1+k^2} &= \dfrac{16}{25} \\   1+k^2 &= \dfrac{25}{16} \\   k^2 &= \dfrac{25}{16} - 1 \\   k^2 &= \dfrac{25}{16} - \dfrac{16}{16} \\   k^2 &= \dfrac{9}{16} \\   k &= \pm\sqrt{\dfrac{9}{16}} \\   k &= \pm\dfrac{3}{4} \end{align*}
Therefore, ignore the negative value (Because \(0<k<1\)), so the final answer is $\color{red}\boxed{\color{black}\dfrac{3}{4}}\color{black}$. We are done ~~ : )

 



\textbf{Solution 7}

Assume that the side length for the largest square in the diagram is $1.$ Since the ratio of the next inner square to the largest square is $k,$ we get that the side length for the second largest square is $k,$ $k^2$ for the third largest square, and $k^3$ for the fourth largest square. So, we can write that 
 \[(1-k^2)+(k^4-k^6)+...=\frac{64}{100}\] by using the diagram since the area of a square with side length $s$ is $s^2.$ This is a geometric sequence with first term $1-k^2$ and ratio $k^4.$ So, we can sum the geometric series for 
 \[\frac{1-k^2}{1-k^4}=\frac{64}{100}.\] Cross multiplying, we have that 
 \[100-100k^2=64-64k^4.\] Moving the terms to one side, we have $64k^4-100k^2+36=0,$ and we can factor as 
 \[(64k^2-36)(k^2-1)=0.\] Since $k<1,$ $k^2-1=0$ is invalid. So, we have 
 \[64k^2-36=0,\] and simplifying we get $k^2=\frac{36}{64}.$ Since $k>0,$ 
 \[k=\frac{6}{8}=\boxed{\frac{3}{4}}.\]




\textbf{Solution 8 (Cheese )}

Measure the side length of the first square to be 2 inches. Measure the side length of the second square to be 1.5 inches. 
 \[\frac{1.5}{2} = \frac{3}{4} = \boxed{\text{(D) } \frac{3}{4}}\]




\textbf{Solution 9 (Also Cheese)}

first square all the answer choices to compare the area. we find that $\frac{16}{25}$  ($\frac{4}{5}$'s square) is too large as it is already 64\% by the first square ring. We find that $\frac{4}{9}$  ($\frac{2}{3}$'s square) is too small as the first ring is 44\% and it drops to ~8\% or so by the 3rd square ring (2nd shaded square ring). This means it has to be bigger than C but smaller than E, and the only answer choice that is that is $\boxed{\text{(D) }\dfrac{3}{4}}$




\textbf{Solution 10 (Fastest and shortest, similar to solution 4, no series)}

Since the figure is self-similar, the nonshaded area is $100\%-64\% = 36\%$ of the largest square and $64\%$ of the square with side length $k$, which has area $k^2$. Thus $36\% = 64\%k^2$, $k^2 = \frac{36}{64}$, so $k = \frac{6}{8} = \boxed{\text{(D) }\dfrac{3}{4}}$.

 



\textbf{Solution 11 (Series Subtraction)}

Assume initial length $1$ since the problem just uses ratios. We see that each side length shrinks by factor $k$, then they are squared to get the areas, which really just multiplies their exponents by $2$. We add our shaded areas and subtract our un-shaded areas, yielding the alternating series 
 \[1 - k^{2} + k^{4} - k^{6}...\]We express this as two infinite series, one being subtracted from the other
 \[\sum_{n=0}^{\infty} k^{4n} -\sum_{n=0}^{\infty} k^{4n+2}\]Rewrite using the infinite geometric series formula and simplify as far as possible

 
 \[\frac{1}{1-k^{4}} - \frac{k^{2}}{1-k^{4}}\]
 \[\frac{1-k^{2}}{1-k^{4}}\]
 \[\frac{1-k^{2}}{(1+k^{2})(1-k^{2})}\]
 \[\frac{1}{1+k^{2}}\]We know this is equal to $\frac{64}{100} = \frac{16}{25}$
 \[\frac{1}{1+k^{2}} = \frac{16}{25}\]
 \[16k^{2}+16 =25\]
 \[16k^{2} = 9\]
 \[k^{2} = \frac{9}{16}\]We must only consider the positive root
 \[k=\boxed{\text{(D) }\dfrac{3}{4}}.\]~shockfront99

 



\textbf{Solution 12 (Quadratics)}

We assign variable $n$ as $\frac{1}{k}.$ For the sake of simplicity, we can assume that the outermost square has a side lenght of $n^2$. Thus, the next two inner side lenghts are $n$ and $1.$  Next, we realize that the square with side lenght $1$ is "similar" to the larger square. Because of this, we know that $64\%$ of it, or $.64$ units squared, is shaded. This means that the shaded area of the larger square is $.64$ units squared combined with the outer shaded ring. Thus:

 
 \[n^4 - n^2 + .64 = .64n^4\]

 To not deal with quartics, we assign $a$ to $n^2.$ Next, we simplify:

 
 \[a^2 - a + .64 = .64a^2\]

 \[100a^2 - 100a + 64 = 64a^2\]

 \[36a^2 - 100a + 64 = 0\]

 \[9a^2 - 25a + 16 = 0\]

 Using the quadratic equation, we obtain values  $\frac{16}{9}$ and $1$. We disregard $1$ because it would cause every square to have the same side lenght. Square rooting $\frac{16}{9}$ yields $\frac{4}{3}$ for $n$ and thus $\frac{3}{4}$ for $k.$



	\item (\textit{AMC 12 12A 2025}) The integer $n$ satisfies the inequalities

 \[\frac{n}{100} < \frac{1}{5 \cdot 10} + \frac{1}{10 \cdot 15} + \frac{1}{15 \cdot 20} + \cdots + \frac{1}{90\cdot 95} + \frac{1}{95\cdot 100} < \frac{n+1}{100}.\]
What is $n$?


 $\textbf{(A)}\   2  \qquad                 \textbf{(B)}\   3  \qquad                 \textbf{(C)}\   4  \qquad                 \textbf{(D)}\   5  \qquad                 \textbf{(E)}\   6$

 \textbf{Answer: }B



\textbf{Solution 1 (20 seconds)}

Pair two students together and put them adjacent on any two seats. There are 6 ways to do this. Considering one of these cases (they are all the same), there are 4 seats left, in which we wish to arrange the teachers together. So pair the teachers together and put them adjacent on any two seats not already occupied by two of the students. There are 3 ways to do this. For all 6 cases, there are \( 6 \times 3 = 18 \) favorable outcomes.

 The number of ways to arrange the 2 students and 2 teachers is $\binom{6}{2} \times \binom{4}{2} = 90$.

 Our probability is \( \frac{18}{90} = \) $\boxed{\textbf{(B) } \frac 15}$




\textbf{Solution 2 - quick (15 seconds)}

We can treat the \(2\) students, \(2\) teachers, and \(2\) empty chairs as \(6\) distinct objects arranged around a circular table. The total number of unique circular arrangements, accounting for rotational symmetry, is:\((6-1)!=5!\)To find the number of favorable arrangements, we group the two students into one object and the two teachers into another object. Combined with the \(2\) distinct empty chairs, we have a total of \(4\) objects to arrange around the circular table.The number of ways to arrange these \(4\) objects in a circle is:\((4-1)!=3!\)We must now account for the internal arrangements within each group. The two students can switch seats in \(2!\) ways, and the two teachers can switch seats in \(2!\) ways.The probability is therefore:\(\frac{3!\cdot 2!\cdot 2!}{5!}=\boxed{\frac{1}{5}}\)~Aar\_v

 



\textbf{Solution 2}

We first count the number of ways to place $4$ distinct people into $6$ distinct chairs: $6\cdot5\cdot4\cdot3 = 360$.

 We now count how many favorable assignments there are. There are $6$ ways to choose an adjacent pair of chairs for the two students. We can arrange the two students in those two chairs in $2$ ways.

 After those two adjacent chairs are taken, there are $4$ chairs left, with $3$ adjacent pairs among them. We choose one of these pairs, arranging the teachers for $3 \cdot 2 = 6$ ways.

 There are $6\cdot 2 \cdot 6 = 72$ favorable arrangements.

 The probability is therefore $\frac{72}{360} = \boxed{\frac{1}{5}}$




\textbf{Solution 3}

We split the problem into cases of where the first teacher and first student sits. The first student sits in a seat with $1$ probability

 


 \textbf{Case $1$:} The first teacher sits next to the first student. There's are $2$ ways to do this so this will happen with $\frac{2}{5}$ chance. Now, there is one valid seat for the second student to sit in with probability $\frac{1}{4}$ and and one valid valid seat for the second teacher to sit in with probability $\frac{1}{3}$. Total probability of this case is $\frac{2}{5}\cdot \frac{1}{4}\cdot \frac{1}{3} = \frac{1}{30}$

 


 \textbf{Case $2$:} The first teacher sits on the opposite end of the circle with $1/5$ chance. This means the second student and the second teacher each have $2$ valid spots to sit in, with probability $\frac{2}{4}$ and $\frac{2}{3}$. Total probability is $\frac{1}{5}\cdot \frac{1}{2}\cdot \frac{2}{3} = \frac{1}{15}$.

 


 \textbf{Case $3$:} The first teacher sits one chair away from the first student with $\frac{2}{5}$ probability because there are $2$ seats that are one chair away. 

 Case $3.1$:

 The second student sits in between the first teacher and first student with $\frac{1}{4}$ chance. The second teacher only has one valid seat, with $\frac{1}{3}$ chance to sit in it. Total probability is $\frac{1}{4}\cdot \frac{1}{3} = \frac{1}{12}$.

 Case $3.2$:

 The second student doesn't sit in between the teacher and the student with a $\frac{1}{4}$ chance. The second teacher can sit in $2$ valid seats with a $\frac{2}{3}$ chance. Total probability is $\frac{1}{4}\cdot \frac{2}{3} = \frac{1}{6}$.

 Total probability for case $3$ is $\frac{2}{5}\cdot (\frac{1}{12} + \frac{1}{6}) = \frac{1}{10}$.

 
Adding up all the cases, you get $\frac{1}{30} + \frac{1}{15} + \frac{1}{10} = \boxed{\frac{1}{5}}$.

 



\textbf{Solution 4}

We consider how many ways there are to place the students and teachers at random at the table. This is a word arrangement in a circle, and so there are $\frac{(6-1)!}{2! \cdot 2!\cdot 2!}=15$ different ways of doing so, assuming that teachers, students and empty spaces are indistinguishable between themselves. 

 Now, we wish to find how many favorable arrangments there are, so suppose we fix two adjacent teachers. Then, there are three "spots" (blocks of two seats) two adjacent students could occupy: two on the side of either teacher, and one such that no student is sitting adjacent to a teacher.

 Thus, the probability is $\frac{3}{15}=\boxed{\frac{1}{5}}$.

 



\textbf{Solution 5 (Step-by-step, needs verification)}

Because this is in a circle, rotating around doesn't increase or decrease the probability (so pin one seat and the others will not rotate).

 Hence, label the spots 1-6 in clockwise. Let student 1 sit in spot 1. Student 2 can therefore sit in either spot 2 or 6, which the probability is 2/5. For explanation purposes, put student 2 in position 2. 

 Here then requires a case breakdown. 

 Case 1: Teacher $1$ sits in spot $3$, hence teacher $2$ needs to sit in spot $4$ to satisfy. This case totals up with $\frac{2}{5}\cdot \frac{1}{4} \cdot \frac{1}{3} = \frac{1}{30}$.

 Case 2: Teacher $1$ sits in spot $4$, hence teacher $2$ needs to sit in spot $3$ or $5$ to satisfy. This case totals up with $\frac{2}{5} \cdot \frac{1}{4} \cdot \frac{2}{3} = \frac{1}{15}$.

 Case 3: Teacher $1$ sits in spot $5$, hence teacher $2$ needs to sit in spot $4$ or $6$ to satisfy. This case totals up with $\frac{2}{5} \cdot \frac{1}{4} \cdot \frac{2}{3} = \frac{1}{15}$.

 Case 4: Teacher $1$ sits in spot $6$, hence teacher $2$ needs to sit in spot $5$ to satisfy. This case totals up with $\frac{2}{5} \cdot \frac{1}{4} \cdot \frac{1}{3} = \frac{1}{30}$.

 By adding the case probabilities, we have $(\frac{1}{30})\cdot 2 + (\frac{1}{15}) \cdot 2 = \boxed{\frac{1}{5}}$.

 Note: Feel free to delete if there is a logic error.~Mitsuihisashi14

 



\textbf{Solution 6 (Quickest)}

For the two teachers' seatings, there are $6$ ways that they are adjacent out of a combination of $n = \binom{6}{2}$ ways. This probability is $\frac{6}{15} = \frac{2}{5}$.For the two students' seatings, there are $3$ ways to choose the two seats out of $n = \binom{4}{2}$ ways. This probability is $\frac{3}{6} = \frac{1}{2}$ . Answer is $\frac{2}{5} \cdot \frac{1}{2} = \boxed{\frac{1}{5}}$.

 



\textbf{Solution 7 (Probability = Desired Arrangements/Total arangements)}

Let the seats be numbered $1,2,3,4,5,6$ in clockwise order. Then assume placing a group of $2$ (either the two students of the two teachers) in seats $1$ and $2$, this group of two can be permuted (or in this case just swpped) in $2!=2$ ways. Now the second group may sit in the seats $3,4$ or $4,5$ or $5,6$ giving a factor of $3$. This group can also be permuted in $2!=2$ ways. This gives 12 ways for the 4 people to be seated under the condition. Without the condition there are $\dfrac{(6-1)!}{2}=60$ ways for the 4 people to be seated (the division by 2! comes from the two empty seats being considered the same). Finally, the desired probability is:
 \[\frac{12}{60}=\boxed{\frac{1}{5}}\]




\textbf{Solution 8}

For the two students, you can seat the first one in any seat, followed by the second one with a $\frac{2}{5}$ possibility of satisfying the condition. After this, we have two cases.

 Case 1: The first teacher sits next to the two students. This can happen $\frac{1}{2}$ of the time and the other teacher has a $\frac{1}{3}$ chance of sitting next to the first one. This gives us a $\frac{1}{2} \cdot \frac{1}{3} = \frac{1}{6}$ chance for this case.

 Case 2: The first teacher sits away from the two students. This can also happen $\frac{1}{2}$ of the time and the other teacher has a $\frac{2}{3}$ chance of sitting next to the first one. This gives us a $\frac{1}{2} \cdot \frac{2}{3} = \frac{1}{3}$ chance for the second case.

 Now, we can find the desired probability as: $(\frac{1}{3} + \frac{1}{6}) \cdot \frac{2}{5} = \frac{1}{2} \cdot \frac{2}{5} = \boxed{\frac{1}{5}}$

 \href{https://y.o.u.t.u.b.e/5Fjos1vBt0A}{https://y.o.u.t.u.b.e/5Fjos1vBt0A}

 


	\item (\textit{AMC 12 12A 2025}) Sharon tosses a fair coin six times and keeps track of how many heads and tails she has gotten. Her first toss is heads. What is the probability that after every toss Sharon's count of heads differs from her count of tails?


 $\textbf{(A)}\   \frac{1}{8}  \qquad                 \textbf{(B)}\   \frac{3}{16}  \qquad                 \textbf{(C)}\   \frac{1}{4}  \qquad                 \textbf{(D)}\   \frac{5}{16}  \qquad                 \textbf{(E)}\   \frac{3}{8}$

 \textbf{Answer: }C



\textbf{Solution 1 (logs and system of equations)}

Substituting $v$ in the equation where $n=5$, we have:

 $\log(k \cdot 5^a m^b)$ $=$ $4+2\log m$.

 Using logarithmic properties, we can write this as:

 $\log(k \cdot 5^a m^b)$ $=$ $\log(10^4 \cdot m^2)$

 We can do this with the other equation where m=25:

 $\log(k \cdot n^a 25^b)$ $=$ $\log(10^4 \cdot n^4)$

 Now we can get rid of the logs on both sides and are left with the following system of equations:

 $k \cdot 5^am^b = 10^4 m^2$
$k \cdot n^a25^b = 10^4 n^4$

 Notice that in the first equation, we can change $m$ arbitrarily, so we know that the exponent of $m$ must be the same - hence $b=2$. Similarly, from the second equation, we get $a=4$. $10^4$ can be written as $2^4 \cdot 5^4$, which means that $k=2^4 = 16$. Thus the answer is $2+4+16 = \boxed{(C)  22}$.

 -Cyrus825 ~ScoutViolet (mostly minor fixes) ~knight10 (minor fixes)

 



\textbf{Solution 2}

We first try to simplify both log equations, and then we bring in the equation for velocity.

 For the equation representing organisms with 5 toes:

 $\log v = 4 + 2 \log m$
$\log v = \log 10^4 + \log m^2$
$\log v = \log(10^4\cdot m^2)$
$v = 10^4\cdot m^2$
$k\cdot 5^am^b = 10^4\cdot m^2$ 

 We do the same with the logarithm equation representing organisms with 25 eyes to get:

 $\log v = 4 + 4 \log n$
$v = 10^4\cdot n^4$
$k\cdot n^a(25^b) = 10^4\cdot n^4$

 Now we need to figure out what $a$, $b$, and $k$ are. Looking at the first equation (or the second), we notice that we can easily match factors. There are two $m$s on the right side so we set $b = 2$ so there are two $m$s on the left side. Now we need $10^4 = k\cdot 5^a = 5^4\cdot 2^4$. There are four $5$s on one side, so we need $a = 4$ to get four $5$s on the other. Now, $k = 2^4 = 16$. This solution works for the first equation.

 Checking if these values work with the second equation, it does, so our answer is $2 + 4 + 16 = \boxed{22}$



	\item (\textit{AMC 12 12A 2025}) A regular octahedron has $8$ equilateral triangle faces with edges of length $\sqrt{2}$. A plane is perpendicular to the line segment joining two opposite vertices of the octahedron and intersects that line segment $\frac{1}{2}$ unit from one of its endpoints. What is the area of the intersection of that plane and the octahedron?


 $\textbf{(A)}\   \frac{1}{4}  \qquad                 \textbf{(B)}\   \frac{1}{2}  \qquad                 \textbf{(C)}\   \frac{1}{2}\sqrt{2}  \qquad                 \textbf{(D)}\   \frac{3}{4}  \qquad                 \textbf{(E)}\   \sqrt{2}$

 \textbf{Answer: }E



\textbf{Solution 1}

We will scale down the diagram by a factor of $3$ so that $AB = 3$ and $AD = 8.$ Since $\angle BEC = 30^{\circ},$ it follows that $\angle BAC = \angle BDC = 30^{\circ}$ as they all subtend the same arc. Similarly, since $\angle CED = 30^{\circ},$ it follows that $\angle CAD = \angle CBD = 30^{\circ}$ as well.

 We obtain the following diagram:
                    \begin{figure}[H]
                        \centering
                        \includegraphics[width=0.5\textwidth]{Solution/AMC_12_12A_2025_8_Solution_Figure_1.png}
                    \end{figure}
                    Note that $\triangle ABD$ has $\angle BAD = 60^{\circ}.$ Applying the law of cosines, we get
 \begin{align*} BD^2 &= AB^2+AD^2-2AB\cdot AD \cdot\cos{60^{\circ}} \\ &= 9 + 64 - 2 \cdot 3 \cdot 8 \cdot \frac{1}{2} \\ &= 49, \end{align*}
from which $BD = 7.$ 

 From here, we wish to find $BF.$ As $AF$ is the angle bisector of $\angle BAD,$ we apply the angle bisector theorem:
 \begin{align*} \frac{AB}{BF} &= \frac{AD}{DF} \\ \frac{3}{BF} &= \frac{8}{7-BF}. \end{align*}
Solving for $BF,$ we get $BF = \frac{21}{11}.$ Remember to scale the figure back up by a factor of $3,$ so our answer is $\frac{21}{11}\cdot 3 = \boxed{\textbf{(E) } \frac{63}{11}}.$

 \textbf{Note} Instead of taking the angle bisectors, one could also take area. The area of $DAB$ is $6\sqrt{3}$. Let $AF = s$ and so through area formulae $2s + \frac{3s}{4} = \frac{11s}{4} = 6\sqrt{3}$ so $s = \frac{24\sqrt{3}}{11}$. Through the law of cosines, $BF = \frac{21}{11}$ so scaled we have $\frac{63}{11}.$




\textbf{Solution 2 Law of (Co)Sine}

From cyclic quadrilateral $CDAE$, we have $\angle CAD = \angle CED = 30^\circ.$ Since $ABDE$ is also cyclic, we have $\angle BAD = \angle BED = 60^\circ$, so, 
 \[\angle BAC= \angle BAD - \angle CAD = 60^\circ - 30^\circ = 30^\circ.\]Using Law of Cosines on $\triangle ABD$, we get 
 \[BD^2=9^2+24^2-2(9)(24)\cos(60^\circ).\] Solving, we get $BD=21$. Next, let $\overline{BF}=x$, and $\angle AFB = \theta$, which means $\overline{FD}=21-x$ and $\angle AFD = 180-\theta$. Using Law of Sines on $\triangle AFB$, we have 
 \[\frac{9}{\sin \theta}=\frac{x}{\sin 30}.\] Solving for $\sin \theta$, we get $\sin \theta = \frac{9}{2x}$. Now we apply the Law of Sines to $\triangle AFD.$ We have 
 \[\frac{24}{\sin(180-\theta)} = \frac{21-x}{\sin 30}.\] Since $\sin(180-\theta) = \sin(\theta),$ and $\sin \theta = \frac{9}{2x}$, we have 
 \[\frac{16x}{3} = 42-2x.\] Solving for $x$ gives $\boxed{x=\frac{63}{11}}$ or $\boxed{\textbf{(E) } \frac{63}{11}}$.

 



\textbf{Solution 3 (Ptolemy’s + Similarity)}

We have $ABCDE$ cyclic, so $\angle BAC=\angle CAD=\angle BEC=30^\circ$. Hence cyclic quadrilateral $ABCD$ has $\angle BAD=60^\circ$. Law of Cosines on triangle $BAD$ gives $\overline{BD}^2=9^2+24^2-2\cdot9\cdot24\cos60^\circ$. Hence $\overline{BD}=21$. Since triangle $BCD$ is a 120-30-30 triangle, we can use law of sines or just memorize ratios to get $\overline{BC}=\overline{CD}=7\sqrt3$. Now Ptolemy’s on $ABCD$ yields $7\sqrt3(9+24)=21\overline{AC}$. Hence $\overline{AC}=11\sqrt3$. Now notice that $\angle BCF=\angle ACB$, and $\angle CBF=\angle CAB=30^\circ$. Hence triangles $CBF$ and $CAB$ are similar, and $\frac{\overline{BF}}{\overline{BC}}=\frac{\overline{AB}}{\overline{AC}}$, so $\frac{\overline{BF}}{7\sqrt3}=\frac9{11\sqrt3}$ and $\overline{BF}=\frac{63}{11}$, or $\boxed{\textbf{(E) } \frac{63}{11}}$.

 



\textbf{Solution 4 (Area & Angle Bisector Length)}

As shown in previous solutions, the angles subtending the same arcs give us $\angle BAC = \angle BEC = 30^\circ$ and $\angle CAD = \angle CED = 30^\circ.$ This tells us that 
 \[\angle BAD = \angle BAC + \angle CAD = 60^\circ,\]and line $AF$ acts as the interior angle bisector of $\angle BAD.$

 We can find the length of the angle bisector $AF$ by splitting the area of $\triangle ABD$ into the sum of the areas of $\triangle ABF$ and $\triangle AFD$:
 \[\text{Area}(\triangle ABD) = \text{Area}(\triangle ABF) + \text{Area}(\triangle AFD).\]

 Using the sine formula for area ($\text{Area} = \frac{1}{2}ab\sin C$), we have:
 \[\frac{1}{2}(AB \cdot AD) \sin 60^\circ = \frac{1}{2}(AB \cdot AF) \sin 30^\circ + \frac{1}{2}(AD \cdot AF) \sin 30^\circ.\]

 Substituting the known values $AB = 9$ and $AD = 24$:
 \[\frac{1}{2}(9 \cdot 24) \frac{\sqrt{3}}{2} = \frac{1}{2}(9 \cdot AF) \frac{1}{2} + \frac{1}{2}(24 \cdot AF) \frac{1}{2}.\]

 Multiplying the entire equation by $4$ to clear the denominators yields:
 \begin{align*} 216\sqrt{3} &= 9(AF) + 24(AF) \\ 216\sqrt{3} &= 33(AF) \\ AF &= \frac{216\sqrt{3}}{33} = \frac{72\sqrt{3}}{11}. \end{align*}


 Now, we apply the Law of Cosines on $\triangle ABF$ to solve directly for $BF$:
 \[BF^2 = AB^2 + AF^2 - 2(AB \cdot AF)\cos 30^\circ.\]

 Substituting $AB = 9,$ $AF = \frac{72\sqrt{3}}{11},$ and $\cos 30^\circ = \frac{\sqrt{3}}{2}$:
 \begin{align*} BF^2 &= 9^2 + \left(\frac{72\sqrt{3}}{11}\right)^2 - 2 \cdot 9 \cdot \frac{72\sqrt{3}}{11} \cdot \frac{\sqrt{3}}{2} \\ BF^2 &= 81 + \frac{15552}{121} - \frac{1944}{11}. \end{align*}


 Finding a common denominator of $121$:
 \begin{align*} BF^2 &= \frac{81 \cdot 121 + 15552 - 1944 \cdot 11}{121} \\ BF^2 &= \frac{9801 + 15552 - 21384}{121} \\ BF^2 &= \frac{3969}{121}. \end{align*}


 Taking the square root of both sides ($\sqrt{3969} = 63$ and $\sqrt{121} = 11$), we find:
 \[BF = \frac{63}{11}.\]

 Thus, the answer is $\boxed{\textbf{(E) } \frac{63}{11}}.$



	\item (\textit{AMC 12 12A 2025}) Suppose that $a$ and $b$ are positive integers satisfying $20a + 26b = 2026$. What is the least possible value of $a+b$?


 $\textbf{(A)}\   71  \qquad                 \textbf{(B)}\   80  \qquad                 \textbf{(C)}\   83  \qquad                 \textbf{(D)}\   86  \qquad                 \textbf{(E)}\   88$

 \textbf{Answer: }E



\textbf{Solution 1 (Rectangular Form)}

We begin by calculating $w^2$:
 \[w^2 = (2+i)^2 = 4+4i-1 = 3+4i.\]Values on the complex plane can easily be represented as points on the Cartesian plane, so we go ahead and do that so we are in a more familiar place. Translating onto the Cartesian plane, we have the points $(2,1)$ and $(3,4)$. The slope of the line passing through these points is $\frac{4-1}{3-2} = 3$, so the equation of this line is
 \begin{align*} y &= 3(x-2)+1 \\ y &= 3x-5. \end{align*}
We want the real number that passes through this line, which is equivalent to the $x-$intercept. This occurs when $y=0$, so the $x$-intercept of this line is $x=\boxed{\textbf{(E)}~\frac53}.$




\textbf{Solution 2 (Polar Form)}

Recall that the slope of a line is $m=\tan\phi,$ where $\phi$ is the angle formed by the line and the positive $x$-axis.

 Note that $|w|=\sqrt{2^2+1^2}=\sqrt{5}.$ In polar coordinates, let $w=\sqrt{5}\operatorname{cis}\theta.$ It follows that $\tan\theta=\frac12.$

 By De Moivre's Theorem, we have $w^2=5\operatorname{cis}(2\theta),$ from which 
 \[\tan(2\theta)=\frac{2\tan\theta}{1-\tan^2\theta}=\frac{4}{3}.\]

 We obtain the following diagram:
                    \begin{figure}[H]
                        \centering
                        \includegraphics[width=0.5\textwidth]{Solution/AMC_12_12A_2025_9_Solution_Figure_1.png}
                    \end{figure}
                    Since $\left|w^2\right|=5,$ we have $w^2=3+4i$ by a $3$-$4$-$5$ triangle. The complex numbers $w$ and $w^2$ correspond to the points $(2,1)$ and $(3,4),$ respectively, and $r$ corresponds to the $x$-intercept of this line. In slope-intercept form, the line containing these two points is $y=3x-5.$ Therefore, the $x$-intercept is $r=\boxed{\textbf{(E)}~\frac53}.$




\textbf{Solution 3 (Coordinate Conversion, Similar to Solution 1)}

Determine $2+i = (2,1)$ on the real cartesian plane and notice $(2+i)^2 = 3+4i$ takes the point $(3,4)$. The slope of the line through the two points is $3x-5 = y$ so $y = 0$ gives $x = \boxed{\frac{5}{3}}$. Because the $x$ axis and the $\Re(x)$ axis are identical, we are finished.

 The solution is here to show a different point of logic for the $x$ axis.

 


	\item (\textit{AMC 12 12A 2025}) Point $C$ lies on a circle with diameter $\overline{AB}$. Point $D$ lies on $\overline{AB}$, $14$ units from $A$, $24$ units from $B$, and $16$ units from $C$. What is the area, in square units, of $\triangle ABC$?


 $\textbf{(A)}\   186\sqrt{2}  \qquad                 \textbf{(B)}\   152\sqrt{3}  \qquad                 \textbf{(C)}\   264  \qquad                 \textbf{(D)}\   187\sqrt{2}  \qquad                 \textbf{(E)}\   108\sqrt{6}$

 \textbf{Answer: }B



\textbf{Solution 1 (Simple)}

The ratio between the radius and the arc length is constant. For the inner circle, the radius, which we will denote as $r$, has a corresponding arc length of $2\pi r - 2\pi$ (the circumference minus the major arc length). For the outer circle, the radius, which is $r + 2\pi$, has a corresponding arc length of $2\pi$. We therefore write the equation 

 
 \[\frac{r}{2\pi r - 2\pi} = \frac{r+2\pi}{2\pi},\]which simplifies to
 \[r = (r+2\pi)(r-1)\]
 \[0 = r^2 + (2\pi-2)r - 2\pi.\]Applying the Quadratic Formula, we get that $r =\boxed{\textbf{(B)} 1 - \pi + \sqrt{1 +\pi^2}}$

 Note: The equation is derived by using the circumference formula and turning that into a ratio. ~AncientKarnivore

 



\textbf{Solution 2}

Let the length of $OA=r$, which is the radius of the smaller circle. Then, the radius of the larger circle, $OB$, is equal to $r+2\pi$. Indeed, we know that the length of major arc $\widehat{AD}=2\pi$ and the length of minor arc $\widehat{BC}=2\pi$. So, using the formula for length of an arc formed by the central angle $\angle COB$, which we denote as $\theta$, we have that:

 
 \[\widehat{AD}=2\pi=(2\pi-\theta) \cdot r,\]
 \[\widehat{BC}=2\pi=\theta \cdot (2\pi+r).\]

 Expanding, we have 
 \[\theta \cdot r +2 \pi \theta = 2\pi,\]
 \[2 \pi r-\theta r=2\pi,\]and by adding the two equations we have that 

 
 \[2 \pi \theta + 2 \pi r =4\pi \implies \theta + r=2.\]

 Indeed, the question is asking for us to solve for $r$, and so we use $\theta = 2-r$ back into our original equation to solve:

 
 \[(2 \pi - \theta) \cdot r = (2 \pi - 2 +r) \cdot r = (2 \pi -2)r +r^2= 2\pi \implies\]
 \[r^2+(2 \pi -2)r - 2 \pi =0.\]

 Using the quadratic formula, we have that 
 \[r=\frac{2-2\pi \pm \sqrt{(2 \pi -2)^2 - 4(1)(-2\pi)}}{2}=\frac{2-2\pi \pm \sqrt{4 \pi^2 -8 \pi +4+8 \pi}}{2}= \frac{2-2\pi \pm 2 \sqrt{1+ \pi^2}}{2}= 1- \pi \pm \sqrt{1+\pi^2}.\]

 Since length must be positive, we consider only the positive root, and so the answer is $\boxed{\textbf{(B)} 1 - \pi + \sqrt{1 +\pi^2}}$.

 


	\item (\textit{AMC 12 12A 2025}) A list contains $2026$ statements, each of which is either true or false. The first statement is "The number of false statements in this list is at least $1001$." The second statement is "The number of false statements in this list is at least $1002$." In general, the $k$th statement in the list is "The number of false statements in this list is at least $k+1000$." How many statements in the list are true?


 $\textbf{(A)}\   342  \qquad                 \textbf{(B)}\   513  \qquad                 \textbf{(C)}\   1013  \qquad                 \textbf{(D)}\   1026  \qquad                 \textbf{(E)}\   1513$

 \textbf{Answer: }A



\textbf{Solution 1 (coordinate geometry)}

As we are given the coordinates of all $3$ vertices, we can compute the slope (or gradient) of each side, and thus the slope of each altitude since it must be perpendicular to the corresponding side. We can therefore determine the Cartesian equation of each altitude using this slope and a point on it (namely, the vertex opposite to the corresponding side), after which we only need to find the point of intersection of at least $2$ of these lines.

 Accordingly, since $BC$ is horizontal (as $B$ and $C$ both have the same $y$-coordinate, namely $27$), we observe that the altitude through $A$ perpendicular to $BC$ will be vertical. As $A$ has $x$-coordinate $2$, the equation of this altitude is therefore precisely $x = 2$. Now, as explained above, we need to determine the equation of $1$ more altitude, so we arbitrarily choose the one through $C$ perpendicular to $AB$. The slope of $AB$ is $\frac{27-31}{8-2} = -\frac{2}{3}$, so because the slopes of perpendicular lines (other than a horizontal and a vertical line) have product $-1$, the slope of this altitude is $\frac{-1}{\left(-\frac{2}{3}\right)} = \frac{3}{2}$. Hence, as this altitude passes through $C$, its equation in point-slope form is $y-27 = \frac{3}{2}(x-18)$, which becomes $y-27 = \frac{3}{2}x-27$, and thus reduces to $y = \frac{3}{2}x$. Consequently, solving this equation simultaneously with $x = 2$ (from above) gives the coordinates of the orthocenter as simply $\left(2,\frac{3}{2} \cdot 2\right) = (2,3)$, and so, finally, the answer is $2+3 = \boxed{\textbf{(A) }5}$.

 \textbf{Alternative: using the altitude through $B$} While we arbitrarily chose to use the altitudes through $A$ and $C$ in the solution above, we could of course use any $2$ of the $3$ altitudes to find the coordinates of the orthocenter as their point of intersection. It's also a good idea to check our answer from above by verifying that $(2,3)$ indeed lies on the third altitude, namely that through $B$ perpendicular to $AC$. Therefore, we compute the slope of $AC$ as $\frac{27-31}{18-2} = -\frac{1}{4}$, so that of the altitude is $\frac{-1}{\left(-\frac{1}{4}\right)} = 4$ (again using the 'slopes have product $-1$' result referred to above). We deduce that the equation of this altitude in point-slope form is $y-27 = 4(x-8) \iff y-27 = 4x-32 \iff y = 4x-5$, and solving this equation simultaneously with either $x = 2$ or $y = \frac{3}{2}x$ from above indeed gives the solution $x = 2, y = 4 \cdot 2 - 5 = \frac{3}{2} \cdot 2 = 3$. In other words, the coordinates of the orthocenter are precisely $(2,3)$, confirming our answer from above.

 



\textbf{Solution 2 (trigonometry)}


                    \begin{figure}[H]
                        \centering
                        \includegraphics[width=0.5\textwidth]{Solution/AMC_12_12A_2025_11_Solution_Figure_1.png}
                    \end{figure}
                    

 Let $H_1$ be the foot of the altitude from $A$ to $BC$, $H_2$ similarly be the foot of the altitude from $B$ to $AC$, and $E$ be the point of intersection of the lines $AH_1$ and $BH_2$. Then we have $\angle EH_1B = \angle EH_2A = 90^{\circ}$, so since angles in a triangle sum to $180^{\circ}$, we obtain 
 \begin{alignat*}{2}\angle EBH\_1 &= 180\textasciicircum{}{\circ} - \angle EH\_1B - \angle BEH\_1 \\ &= 180\textasciicircum{}{\circ} - 90\textasciicircum{}{\circ} - \angle BEH\_1 \qquad &&\text{(from above)} \\ &= 180\textasciicircum{}{\circ} - \angle EH\_2A - \angle BEH\_1 \qquad &&\text{(again from above)} \\ &= 180\textasciicircum{}{\circ} - \angle EH\_2A - \angle AEH\_2 \qquad &&\text{(as } H\_1 \text{ lies on } AE \text{ and } B \text { lies on } H\_2E \text{, so } \angle BEH\_1 \text{ and } \angle AEH\_2 \text{ are the same angle)} \\ &= \angle EAH\_2.\end{alignat*} Moreover, $\angle EAH_2$ is also the same angle as $\angle H_1AC$, since (once again) $H_1$ lies on $AE$ and $H_2$ lies on $AC$. As $H_1$ has the same $x$-coordinate as $A$ and the same $y$-coordinate as $C$, its coordinates are precisely $(2,27)$, so we now deduce

 
 \begin{alignat*}{2}\tan{\angle EBH\_1} = \tan{\angle EAH\_2} &= \tan{\angle H\_1AC} \\ &= \frac{CH\_1}{AH\_1} \qquad &&\text{(as } \angle AH\_1C = 90\textasciicircum{}{\circ} \text{, i.e. triangle } \triangle AH\_1C \text{ is right-angled at } H\_1 \text{)} \\ &= \frac{18-2}{31-27} \qquad &&\text{(using the coordinates of } H\_1 \text{ from above)} \\ &= 4.\end{alignat*}

 But again as $\angle EH_1B = 90^{\circ}$, triangle $\triangle EH_1B$ is also right-angled at $H_1$, so, similarly to above,

 
 \begin{align*}&\tan{\angle EBH_1} = \frac{EH_1}{BH_1} \\ &\phantom{\tan{\angle EBH_1} \:} = \frac{EH_1}{8-2} \qquad \text{(again using the coordinates of } H_1 \text{ from above)} \\ &\iff \frac{EH_1}{6} \: = 4 \iff EH_1 = 24.\end{align*}


 Hence, as $E$ lies vertically \textit{below} $H_1$, it finally follows that its coordinates are $(2,27-EH_1) = (2,27-24) = (2,3)$, and thus the answer is $2+3 = \boxed{\textbf{(A) }5}$.

 



\textbf{Solution 3 (Graph analysis)}

WLOG assume $BC$ is base as $BC$ is collinear on the same line $y = 27$. Our triangle looks as such:

 
                    \begin{figure}[H]
                        \centering
                        \includegraphics[width=0.5\textwidth]{Solution/AMC_12_12A_2025_11_Solution_Figure_2.png}
                    \end{figure}
                    

 One altitude has slope undefined and equation $x = 2$. Another such altitude defined $h_B$ is perpendicular to $AC$ who has slope $\frac{-1}{4}$, so $h_B$ has slope 4. It passes through $(8,27)$, so $y = 4x+b$ and $27 = 32 + b$ so $b = -5$ and therefore $y = 4x-5$ and when $x = 2$, $y = 8-5 = 3$. The sum is $\boxed{5}$.

 
 \[\textbf{OR}\]

 Translate point $B$ to $(0,0)$ and so $A(-6,4), C(10,0)$. The line $x = -6$ and $y = 4x$ exists, whose intersection is at $(-6,-24)$. Adding back $(x,y) \rightarrow (x+8, y+27)$ we obtain $H(2,3)$ whose sum is $\boxed{5}$.

 



\textbf{Solution 4 (Finding Slope among other Things)}

The orthocenter of a triangle is located at the intersection of its altitudes. Altitudes are lines that cross from a vertex to the opposite side, such that they form a $90^\circ$ angle where they cross the base. We can use this fact and our knowledge of slope to find where two of these lines intersect.

 Start by looking at line $BC$, a horizontal line with a slope of $0$. A line perpendicular to this that crosses point $A$ must have the equation $x=2$.

 We now do the same for line $AB$. The line passing through $AB$ has a slope of\begin{align*}m &= \frac{27-31}{8-2} \Longrightarrow -\frac{2}{3}.\end{align*}
The line perpendicular to this has a slope of $\frac{3}{2}$. This altitude will also cross point $C$. Using point-slope form, and substituting in $x=2$ from our first altitude, we find the $y$-coordinate of the intersection:\begin{align*}y &= \frac{3}{2}(x-18)+27 \Longrightarrow \frac{3}{2}(2-18)+27 = -24+27 = 3.\end{align*}
Thus, the orthocenter is located at $(2, 3)$. The sum of the coordinates is\begin{align*}2+3 &= \boxed{5}\end{align*}
\begin{align*}5 = \boxed{A}\end{align*}


 



\textbf{Solution 5 (Vectors)}

The orthocenter of a triangle is the intersection of any 2 of the altitudes. This is helpful because we can choose any two altitudes (some may be more convenient than others).

 
                    \begin{figure}[H]
                        \centering
                        \includegraphics[width=0.5\textwidth]{Solution/AMC_12_12A_2025_11_Solution_Figure_2.png}
                    \end{figure}
                    

 Diagram by ~pinotation

 Notice it is convenient to use the altitude drawn from vertex $A$, because the line perpendicular to side $BC$ is simply parallel to the y-axis.

 The line parallel to the y-axis that crosses vertex $A$ is simply $x=2$. For the next altitude, we continue with vectors

 
Let $(x, y)$ be point $D$ such that $\overleftrightarrow{AB} \perp \overleftrightarrow{CD}$

 \begin{array}{l}\implies \overrightarrow{AB} \cdot \overrightarrow{CD} = 0 \\\implies ((8 - 2)\hspace{5mm}(27 - 31)) \cdot ((x - 18)\hspace{5mm}(y - 27)) = 0 \\\implies 6(x - 18) + (-4)(y - 27) = 0 \\\implies 4y = 6x \\\implies y = \frac{3x}{2}\end{array}

 By naming connecting a point $D$ to vertex $C$ and defining it as perpendicular to line $\overleftrightarrow{AB}$, we know line $\overline{CD}$ lies on the line containing the altitude from $C$. But because they are perpendicular, the dot product between the two vectors is 0, from which we get the line point $D$ lies on (and therefore the altitude).

 Now finding the intersection of the 2 lines, 

 $x=2$
$y = \frac{3}{2}x$
$\implies$ Orthocenter is  $(2, 3)$

 The answer is $2 + 3 = 5$ $\implies \boxed{A}$



	\item (\textit{AMC 12 12A 2025}) Polly and Quentin are dancing. They start $1$ meter apart. Polly rotates $120^{\circ}$ clockwise around Quentin while he remains still, staying $1$ meter away throughout. Then Quentin rotates $120^{\circ}$ clockwise around Polly while she remains still, staying $1$ meter away throughout. They repeatedly do this pair of moves from their new positions. How many meters from where she started is Polly after they have each moved $10$ times?


 $\textbf{(A)}\   0  \qquad                 \textbf{(B)}\   \frac{\sqrt{3}}{2}  \qquad                 \textbf{(C)}\   1  \qquad                 \textbf{(D)}\   \sqrt{3}  \qquad                 \textbf{(E)}\   2$

 \textbf{Answer: }B



\textbf{Solution 1}

We will need to determine the sum of the reciprocals of the roots. To find the sum of the reciprocals of the roots $p,q$ of the quadratic $ax^2+bx+c$, we use Vieta's formulas. Recall that $p+q = -b/a$ and $pq = c/a$. Therefore, 
 \[\frac{1}{p} + \frac{1}{q} = \frac{p+q}{pq} = \frac{\frac{-b}{a}}{\frac{c}{a}} = \frac{-b}{a} \cdot \frac{a}{c} = \frac{-b}{c},\]which doesn't depend on $a$.

 The sum of the reciprocals of the roots of the quadratic $x^2-4x-3$ is $\frac{-(-4)}{-3} = -4/3.$ The same is true for every quadratic in the form $ax^2-4x-3$. The sum of all the reciprocals of the roots of $ax^2+bx+c$ is $2025 \cdot \left(-\frac{4}{3}\right).$

 Because we have $2025$ quadratics, there are $2 \cdot 2025 = 4050$ total roots. Our answer is $\frac{1}{\frac{1}{4050}\cdot \frac{-4\cdot 2025}{3}} =\boxed{(B)\ -\dfrac{3}{2}}$.

 



\textbf{Solution 2 (similar to solution 1 but with quadratic formula)}

We first find the general roots for quadratics in the form $kx^2-4x-3$. Using the quadratic formula we have \begin{align*}x&=\frac{4 \pm \sqrt{16+12k}}{2k} \\&= \frac{4\pm 2\sqrt{4+3k}}{2k} \\&= \frac{2+\sqrt{4+3k}}{k}, \frac{2-\sqrt{4+3k}}{k} \\\end{align*}
Since we are asked to add the reciprocals of all $4050$ roots in the harmonic mean, we will first add the general roots in terms of $k$. We have,\begin{align*}\frac{k}{2+\sqrt{4+3k}} + \frac{k}{2-\sqrt{4+3k}} &= \frac{k(2-\sqrt{4+3k}) + k(2+\sqrt{4+3k})}{2^2-(\sqrt{4+3k})^2} \\&= -\frac{4k}{3k}=-\frac{4}{3}. \\\end{align*}
Thus, each pair of roots add up to $-\frac{4}{3}$, and since there are $2025$ pairs of roots, the harmonic mean of the desired expression is \begin{align*}\frac{1}{\frac{1}{4050} \left (2025 \left (-\frac43 \right ) \right)} &= \frac{1}{\frac12 \left ( -\frac43 \right)} \\& = \boxed{-\frac32}, \boxed{B}. \\\end{align*}


 



\textbf{Solution 3 (Vieta's only)}

We are asked to find $\frac{4050}{\sum_{i=1}^{4050}\frac{1}{i_1}}$. By Vieta's, note that $r_1 \cdot r_2 \dots r_{4050} = -\frac{3^{2025}}{2025!} = k$ ($k$ is a constant). Then, note that we are asked to find $\frac{4050}{\sum_{i=1}^n \frac{\frac{k}{r_i}}{k}}$, and by Vieta's we get that $\sum_{i=1}^n \frac{k}{r_i} = \frac{2025 \cdot 4 \cdot 3^{2024}}{2025!}$, so substituting this in, we get 
 \[-\frac{4050}{\frac{\frac{2025 \cdot 4 \cdot 3^{2024}}{2025!}}{\frac{3^{2025}}{2025!}}} = -\frac{4050 \cdot 3^{2025}}{4050 \cdot 3^{2024} \cdot 2} = \boxed{-\frac{3}{2}}\].

 


 



\textbf{Solution 4 (Vieta's Formulas)}

We have: 
 \[\prod_{k = 1}^{2025} (kx^{2} - 4x - 3) = (x^{2} - 4x - 3)(2x^{2} - 4x - 3)(3x^{2} - 4x - 3) \cdots (2025x^{2} - 4x - 3)\]

 So we can analyze each part of this "snake":

 (definitions added for clarity in grammar & notation)

 if 
 \[\prod_{k = 1}^{2025} (kx^{2} - 4x - 3)=\prod_{k = 1}^{2025} P_k(x)\]and
 \[P_k(x)= (kx^{2} - 4x - 3)\]

 According to Vieta's Formulas:

 For $k=1$, $P_1(x)=x^{2} - 4x - 3$:We have:
 \[\begin{cases}  x_1+x_2=\frac{4}{1}\\ x_1x_2=-\frac{3}{1} \end{cases}\]So we have: 
 \[\dfrac{1}{x_1}+\dfrac{1}{x_2}=\dfrac{x_1+x_2}{x_1x_2}=-\dfrac{4}{3}\]

 For $k=2$: $P_2(x) = 2x^{2} - 4x - 3$:We have:
 \[\begin{cases}  x_3+x_4=\frac{4}{2}\\ x_3x_4=-\frac{3}{2} \end{cases}\]So we have: 
 \[\dfrac{1}{x_3}+\dfrac{1}{x_4}=\dfrac{x_3+x_4}{x_3x_4}=-\dfrac{4}{3}\]

 For $k=3$, $P_3(x)=3x^{2} - 4x - 3$:We have:
 \[\begin{cases}  x_5+x_6=\frac{4}{3}\\ x_5x_6=-\frac{3}{3} \end{cases}\]So we have: 
 \[\dfrac{1}{x_5}+\dfrac{1}{x_6}=\dfrac{x_5+x_6}{x_5x_6}=-\dfrac{4}{3}\]

 ......

 Then we should check the last k, For $k=2025$, $P_{2025}(x) = 2025x^{2} - 4x - 3$:We have:
 \[\begin{cases}  x_{4049}+x_{4050}=\frac{4}{2025}\\ x_{4049}x_{4050}=-\frac{3}{2025} \end{cases}\]So we have: 
 \[\dfrac{1}{x_{4049}}+\dfrac{1}{x_{4050}}=\dfrac{x_{4049}+x_{4050}}{x_{4049}x_{4050}}=-\dfrac{4}{3}\]
 \begin{align*} \therefore H.M.&=\dfrac{1}{\dfrac{1}{4050}\sum_{i=1}^{4050}\dfrac{1}{x_i}}\\ &=\dfrac{4050}{2025\times\left(-\dfrac{4}{3}\right)}\\ &=-\dfrac{3}{2} \end{align*}
Therefore, the final answer is $\color{red}\boxed{\color{black}-\dfrac{3}{2}}\color{black}$.

 



\textbf{Solution 5 (Vieta's Formulas)}

For a general quadratic polynomial $ax^2 + bx + c = 0$ with roots $p$ and $q$, the sum of the reciprocals is:\begin{align*}\frac{1}{p} + \frac{1}{q} &= \frac{p+q}{pq} \\\end{align*}
Using Vieta's formulas ($p+q = -\frac{b}{a}$ and $pq = \frac{c}{a}$):\begin{align*}\frac{1}{p} + \frac{1}{q} &= \frac{-\frac{b}{a}}{\frac{c}{a}} = -\frac{b}{c}\end{align*}
This result demonstrates that the sum of the reciprocals depends only on the coefficients $b$ and $c$, and is independent of the leading coefficient $a$. Since the overall polynomial is a product of quadratics, and all these quadratics share the same $b$ and $c$ coefficients, the sum of the reciprocals for the roots of every individual quadratic is the same. We use the coefficients from one of the quadratics, which is $ax^2 - 4x - 3 = 0$. Thus, $b=-4$ and $c=-3$.

 
 \[\sum (\text{reciprocals}) = -\frac{b}{c} = -\frac{(-4)}{(-3)} = -\frac{4}{3}\]

 We substitute the calculated sum of the reciprocals into the formula for the harmonic mean of 2 numbers:

 
 \[\text{HM} = \frac{2}{-\frac{4}{3}} = 2 \times \left(-\frac{3}{4}\right) = -\frac{6}{4} = -\frac{3}{2}\]

 The harmonic mean of the roots of any single quadratic is $-\frac{3}{2}$. Since the sum of the reciprocals is the same for all quadratics, the harmonic mean of all the roots of the overall polynomial is also $\boxed{\textbf{(B) }\dfrac{-3}{2}}$.

 \textbf{Note} It is important to note that the question asks for the sum of all \textit{real} roots. We must therefore be careful to make sure that the $4050$ roots of the polynomial are all real, and also that they are distinct, as otherwise e.g. double roots would be counted twice, when the problem requires us to only count each root once (i.e. regardless of multiplicity). Accordingly, we compute that each quadratic factor $\left(kx^2-4x-3\right)$ (for $1 \leq k \leq 2025$) has discriminant $(-4)^2-4 \cdot k \cdot (-3) = 16+12k$, which is clearly positive as $k \geq 1$. This means the $2$ roots of each quadratic factor are indeed real and distinct (i.e. none of these quadratics has a double root), so it only remains to show that these roots are also distinct from those of all the \textit{other} quadratic factors. Hence, if we assume that $a$ is a common root of $2$ different quadratic factors simultaneously, say $\left(px^2-4x-3\right)$ and $\left(qx^2-4x-3\right)$ with $1 \leq p,q \leq 2025$ and $p \neq q$, then we must have $pa^2-4a-3=qa^2-4a-3=0$, which implies $pa^2 = qa^2$. This equation factorises as $(p-q)a^2 = 0$, and because $p \neq q$, i.e. $p-q \neq 0$, it further reduces to simply $a^2 = 0$, and hence $a = 0$. But then it follows that $pa^2-4a-3 = p \cdot 0^2 - 4 \cdot 0 - 3 = -3 \neq 0$, yielding a contradiction, so $a$ cannot in fact have been a common root of $2$ different quadratic factors, and thus we finally deduce that all $4050$ roots are distinct, as required.~ Shadowleafy

 



\textbf{Solution 6}

Let the polynomial be $f(x),$ and denote the $4050$ roots to be $x_1,x_2,...,x_{4050}.$ Hence, 
 \[HM = \dfrac{4050}{\frac{1}{x_1}+\frac{1}{x_2}+...\frac{1}{x_{4050}}}.\]We can multiply the numerator and denominator of this fraction by $x_1x_2...x_{4050}$ to create symmetric sums, which yields 
 \[HM = \dfrac{4050(x_1x_2...x_{4050})}{x_2x_3...x_{4050}+x_1x_3...x_{4050}+...+x_1x_2...x_{4049}}.\]

 By Vieta's Formulas, since $f(x)$ is of even degree, the product of its roots, $x_1x_2...x_{4050},$ is just the constant term of $f(x),$ call it $c_0.$ Likewise, the denominator of our harmonic mean, $x_2x_3...x_{4050}+x_1x_3...x_{4050}+...+x_1x_2...x_{4049},$ is the negated coefficient of $x$ in the standard form of $f(x).$ Let the coefficient of $x$ in the standard form of $f(x)$ be $c_1.$ Note that we do not have to worry about dividing by the coefficient of $x^{4050}$ when using Vieta's Formulas, as they will eventually cancel out in our harmonic mean calculations.

 So, 
 \[HM = \dfrac{4050c_0}{-c_1}.\]

 The constant term in $f(x)$ is just $c_0=(-3)^{2025}.$ For the coefficient of the $x$ term in $f(x),$ there are $\dbinom{2025}{2024}=2025$ ways to choose $2024$ of the trinomials to include a $-3,$ and the one trinomial not chosen will include a $-4x.$ Hence, $c_1=2025\cdot (-3)^{2024}\cdot (-4).$

 Finally, 
 \[HM=\dfrac{4050\cdot(-3)^{2025}}{-2025\cdot (-3)^{2024}\cdot (-4)}=\dfrac{2\cdot(-3)}{-(-4)}=\boxed{\text{(B) }-\dfrac{3}{2}}.\]




\textbf{Solution 7 (Cheese )}

We can try replacing the $2025$ in the problem with smaller numbers like $1$ and $2$. For $1$, we have the quadratic $x^2-4x-3$. Using the quadratic formula, we have the roots are $2+\sqrt{7}$ and $2-\sqrt{7}$ The harmonic mean of these two is 
 \[\frac{2}{\frac{1}{2+\sqrt{7}} + \frac{1}{2-\sqrt{7}}}=\frac{2}{\frac{\sqrt{7}-2}{3}+\frac{2+\sqrt{7}}{-3}}=\frac{2}{\frac{-4}{3}}=-\frac{3}{2}\].

 We notice that this is one of the answer choices. Also, given that the random choice of $2025$, and that the rest of the answer choices are also simple fractions, we can safely guess that the answer is $\boxed{\textbf{(B)}~-\frac{3}{2}}$




\textbf{Solution 8 (Symmetry of Unweighted Means)}

Note that all the roots are real since for each quadratic, 
 \[\Delta=b^2-4ac=(-4)^2-4k(-3)=16+12k>0\] for $k\ge 1$. We will use the fundamental theorem of symmetric polynomials. By Vieta's formulas, we have for each quadratic $ax^2+bx+c=kx^2-4x-3$:
 \[\frac{1}{r}+\frac{1}{s}=\frac{r+s}{rs}=\frac{-b/a}{c/a}=-\frac{b}{c}=-\frac{4}{3}\]Thus, $\mathrm{HM}(r,s)=\frac{2}{\frac{1}{r}+\frac{1}{s}}=\frac{2}{-\frac{4}{3}}=-\frac{3}{2}$. Since the unweighted harmonic mean is symmetric in its arguments, we can group terms to get
 \[\mathrm{HM}(r_{1},r_{2},r_{3},\dots,r_{4050})=\mathrm{HM}\bigg(\mathrm{HM}(r_{1},r_{2}),\mathrm{HM}(r_{3},r_{4}),\dots,\mathrm{HM(r_{4049},r_{4050})}\bigg)=\boxed{\textbf{(B)}~-\frac{3}{2}}\]

 Note: For an unweighted symmetric mean $M$, we have for $d|n$, $\mathrm{M}\{ x_{k} \}_{k=1}^{n}=\mathrm{M}\{ \mathrm{M}\{ x_{ds+k} \}_{k=1}^{d} \}_{s=0}^{n/d-1}$ and $\mathrm{M}\{ \bar{x} \}_{k}=\bar{x}$ for constant mean $\bar{x}$. You can also check this individually for HM: 
 \[\mathrm{HM}(x_{1},\dots,x_{n})=n(x_{1}^{-1}+\cdots+x_{n}^{-1})^{-1}=\frac{n}{d}\left[ \underbrace{ d\left(x_{1}^{-1}+\cdots+x_{d}^{-1}\right)+\cdots+d\left(x_{(n-1)d+1}^{-1}+\cdots+x_{n}^{-1}\right) }_{ \frac{n}{d}\text{ terms} } \right]\]or by checking arithmetic mean and using $\mathrm{HM}\{ x_{k} \}_{k}=\frac{1}{\mathrm{AM}\left\{ \frac{1}{x_{k}} \right\}_{k}}$.

 



\textbf{Solution 9 (Easy-ish to understand)}

Note: I will try my best to make this solution look less intimidating to read as the other ones. This is my third time writing a solution, so I'm not very used to the $\LaTeX$. 

 The harmonic mean that the question is asking is 
 \[\frac{1}{\frac{r_1+r_2+...r_{4049}+r_{4050}}{4050}}\] where $r_i$ is between $1$ and $4050$ and is the reciprocal of all $4050$ roots of the polynomial. All we need to find is the sum of the reciprocals of the roots and we can substitute that into this formula. 

 First of all, notice that there are $2025$ quadratic polynomials. They each hold $2$ out of the $4050$ roots, and the sum of the reciprocals of their roots are $\frac{1}{a} + \frac{1}{b} = \frac{a+b}{ab}.$ 

 Using Vieta's we get $\frac{a+b}{ab}=-\frac{4}{3}$ on the every quadratic because only the coefficient for $x^2$ changes. Also, $a+b$ and $ab$ have the denominator of $a$, so when they are divided, the quotient won't change. Therefore, the sums of the reciprocals of the quadratics sum to $2025 \cdot \frac{-4}{3}$.

 Now we can substitute: 
 \[\frac{1}{\frac{2025 \cdot \frac{-4}{3}}{4050}}.\] This gets us \begin{align}\frac{1}{\frac{2025 \cdot \frac{-4}{3}}{4050}} &= \frac{1}{\frac{2025 \cdot (-4)}{3 \cdot 4050}} \\&= \frac{1}{\frac{-8100}{12150}} \\&= \frac{1}{\frac{-2}{3}} \\&= 1 \cdot \frac{3}{-2} \\&= -\frac{3}{2}\end{align}

 Or...you could just solve the last part yourself easier than my solution!

 



\textbf{Solution 10 (reciprocal roots)}

We know that if for polynomial $ax^2 + bx + c$ is $r_1$ and $r_2$, then the polynomial with roots $\frac{1}{r_1}$ and $\frac{1}{r_2}$ (reciprocal roots) will be $cx^2 + bx + a$ where the coefficients are flipped (works for any degree polynomial, not just quadratic). Thus, using vieta's formula, the sum of the reciprocal roots for polynomial $kx^2 - 4x - 3$ will be the sum of the roots for polynomial $-3x^2 - 4x + k$ which is $- \frac{-4}{-3}=- \frac{4}{3}$.

 There are $2025$ polynomials like this, so the sum of roots would be $2025\cdot (- \frac{4}{3})$. Finding the harmonic mean, we get:

 
 \[\frac{1}{\frac{1}{4050} \cdot \frac{-4\cdot 2025}{3}} = \frac{1}{\frac{-2}{3}}\]

 Or $\boxed{\textbf{(B) }-\frac{3}{2}}$



	\item (\textit{AMC 12 12A 2025}) A jar contains $12$ green marbles and $8$ red marbles. In a game, Anika and Ben take turns drawing a random marble from the jar, with replacement. If Ben ever draws a red marble, he wins, and if Anika ever draws a green marble, she wins. Otherwise, the game continues with the next player's turn. The probability that Ben eventually wins is $A$ if Anika goes first and $B$ if Ben goes first. What is $B - A$?


 $\textbf{(A)}\   \frac{2}{25}  \qquad                 \textbf{(B)}\   \frac{4}{25}  \qquad                 \textbf{(C)}\   \frac{4}{19}  \qquad                 \textbf{(D)}\   \frac{6}{19}  \qquad                 \textbf{(E)}\   \frac{9}{20}$

 \textbf{Answer: }D



\textbf{Solution 1}

We first find what $N$ is by figuring out how many numbers we need to take out of the set so that the set does not contain $5$ consecutive integers. Since $N$ must be maximized, we must minimize what numbers are removed, and we quickly find that taking two numbers out works. Consider taking out $5$ and $10$. You are left with $\{1,2,3,4,6,7,8,9,11,12,13\}$, which does not have a string of $5$ consecutive integers.

 There are only $3$ ways to take out two integers such that the resulting set meets our condition ($5$ and $10$, $5$ and $9$, or $4$ and $9$), and ${\binom{13}{2}}=78$ total ways to choose such integers. Therefore, the probability is $\boxed{\dfrac{1}{26}}$.

 Minor edits ~aashrithm29

 Very minor LaTeX edit ~ PerseverePlayer

 Note: An alternative way to come to the conclusion that there are only three possible subsets of length twelve is as follows:

 Let the two numbers removed from $C$ be $a$ and $b$ such that $a<b$.  Then, because we cannot have any sequence of five consecutive numbers, $a \leq 5$, $b-a \leq 5$, and $13 - b +1\leq 5$.  Then, graphing this system in the $ab$ plane, we easily see that there are only three solutions for $(a,b)$: $(5,10), (4,9)$, and $(5,9)$. Then, simply proceed as described above. 

 



\textbf{Solution 2}

Trying to find a subset that satisfies the condition, we get $\{1,2,3,4,6,7,8,9,11,12,13\}$, which has $N=11$ elements. The subsets $\{1,2,3,5,6,7,8,10,11,12,13\}$ and $\{1,2,3,4,6,7,8,10,11,12,13\}$ also work. In total, we have $3$ subsets and $\binom{13}{11}$ ways to choose $11$ elements from $C$, so the probability is $\dfrac{3}{\binom{13}{11}} = \dfrac{1}{26}$. Thus, the answer is $\boxed{\textbf{(D)}}$~anzhuPro

 



\textbf{Solution 3}

First, let's find the number $N$:Consider if:

 
$N=13$: 

 It is impossible - It doesn't remove any integer, violating the condition "not contain five consecutive integers".

 
$N=12$:

 It is also impossible, and still does not satisfy the condition. An example could be $\{1,2,3,4,6,7,8,9,10,11,12,13\}$. $1$, $2$, $3$, and $4$ doesn't violate the condition, but $6-13$ are all consecutive, hence $N$ can't be $12$. 

 
$N=11$: 

 This is possible, because $\dfrac{11}{3} = 3$ remainder $2$. This means that we can have one string of consecutive length $3$ and two strings of consecutive length $4$. Hence, it could be either:

 
            \begin{itemize}
        \item $\{1,2,3,5,6,7,8,10,11,12,13\}$ with consecutive integer lengths of $(3, 4, 4)$.
        \item $\{1,2,3,4,6,7,8,10,11,12,13\}$ with consecutive integer lengths of $(4, 3, 4)$.
        \item $\{1,2,3,4,6,7,8,9,11,12,13\}$ with consecutive integer lengths of $(4, 4, 3)$.
            \end{itemize}
            Thus giving $3$ cases. The total ways of picking out $2$ integers from a set of $13$ is $\binom{13}{2} = 78$. Hence, the probability is $\dfrac{3}{78} = \boxed{\dfrac{1}{26}}$.

 
Note: You can also consider the combinations of 1 string of consecutive length 3 and 2 strings of consecutive length 4. The length of 3 has 3 places to go, filling the rest 2 with 4 places each, giving 3/78 as well, resulting in the same answer. 

 Formatting and minor edits ~Logibyte

 


	\item (\textit{AMC 12 12A 2025}) A line segment is drawn between two points on the graph of $y = x^{2}$ whose $x$-coordinates have been chosen independently and uniformly at random in the interval $[0,2]$. The probability that the $y$-coordinate of the midpoint of the segment is greater than $2$ lies in which of the following intervals?


 $\textbf{(A)}\   (0.20, 0.21)  \qquad                 \textbf{(B)}\   (0.21, 0.22)  \qquad                 \textbf{(C)}\   (0.22, 0.23)  \qquad                 \textbf{(D)}\   (0.23, 0.24)  \qquad                 \textbf{(E)}\   (0.24, 0.25)$

 \textbf{Answer: }D



\textbf{Solution 1 (Simple)}

Let the outer ellipse be ellipse 1, and the inner ellipse be ellipse 2. Let $a_1$, $b_1$, and $c_1$, correspond to the semimajor axis, semiminor axis, and focal distance of ellipse $1$, respectively. Similarly, let $a_2$, $b_2$, and $c_2$ correspond to the semimajor axis, semiminor axis, and focal distance of ellipse $2$, respectively.

 Ellipses with the same eccentricity are similar, so $\frac{a_1}{a_2} = \sqrt{2025} = 45$ because the ratio of semimajor axes between similar ellipses is equal to the square root of the ratio between their areas. Notice how
 \[a_1 = c_1 + c_2 + a_2.\]Substituting from the eccentricity equation,
 \[a_1 = ea_1 + ea_2 + a_2.\]Rearranging gives us
 \[a_1 - ea_1 = ea_2 + a_2\]
 \[a_1(1-e) = a_2(1+e)\]
 \[\frac{a_1}{a_2} = \frac{1+e}{1-e} = 45.\]Solving for $e$ then yields $\boxed{\text{(D) }\dfrac{22}{23}}.$

 \textbf{Conic Sections Solution} 
                    \begin{figure}[H]
                        \centering
                        \includegraphics[width=0.5\textwidth]{Solution/AMC_12_12A_2025_14_Solution_Figure_1.png}
                    \end{figure}
                    Let $A$ and $B$ be the centers of the ellipses with foci $F, G$ and foci $G, H$ respectively. Then, say $AQ=p$, $AP = q$, $BS=r$, and $BR=s$. Through area formula of ellipse ($ab\pi$, where $a$ is half the major axis and $b$ is half the minor axis), the area of the ellipse with center at $A$ is $\pi pq$ and with center at $B$ is $\pi rs$. Then, we have $\frac{\pi pq}{\pi rs} = 2025$, in which $\frac{pq}{rs} = 2025$.

 Additionally, when two ellipses have the same eccentricity, that essentially means the ellipses are stretched the same amount. Therefore, $\frac{p}{q} = \frac{r}{s}$, in which $ps = qr$, and $p=\frac{qr}{s}$.

 From the previous equation $\frac{pq}{rs}=2025$, we have $pq=2025rs$, in which $\frac{q^2r}{s} = 2025rs$, in which $q^2r = 2025rs^2$, in which $q^2 = 2025s^2$, and $q=45s$ ( $\neg -45$ as length cannot be negative). Similarly, $p = 45r$.

 
                    \begin{figure}[H]
                        \centering
                        \includegraphics[width=0.5\textwidth]{Solution/AMC_12_12A_2025_14_Solution_Figure_2.png}
                    \end{figure}
                    Connect $F$ to $P$ and $G$ to $R$. We form two similar triangles $PAF$ and $RBG$ (through ASA). We just need to determine the length $BG$, and the problem is solved. Notice that $\frac{RB}{PA} = \frac{BG}{AF}$. because $RB:PA = 45:1$, $BG:AF = 1:45$, and if $AF = 45x$, $BG = x$. To solve for $x$, we connect $B$ to the tangency point of both ellipses ($W$). Then, $AW = 45r$ and $BW = r$, so $AB=44r$. We found that $AF = 45x$, and through definition of foci, $AG = 45x$. $GB = x$, so $AB=46x$. However, $AB = 44r$, so $46x=44r$, $x = \frac{22r}{23}$.

 The length of the major axis is $2r$, and because eccentricity is $\frac{c}{a}$, where in our case, $c=x$ and $a=r$, we have $\frac{22r}{23r}$, or $\boxed{\frac{22}{23}}.$




\textbf{Solution 2}

We label the outer ellipse as Ellipse $1$ and the inner ellipse as Ellipse $2$. Let the focal distance of Ellipse $1$ be $2c_1$ and the major axis be $2a_1$. Similarly, Ellipse $2$ has a focal distance of $2c_2$ and a major axis of $2a_2$. 

 We label the rightmost vertex of both ellipses as $I$. Because $I$ lies on Ellipse $1$, the sum of the lengths of $FI$ and $GI$ equals $2a_1$. The length of $FI$ can be written as $c_2+a_2+2c_1$, and the length of $GI$ can be written as $c_2+a_2$. Therefore, 
 \[(c_2+a_2+2c_1) + (c_2+a_2) = 2a_1\]
 \[2c_2+2a_2+2c_1 = 2a_1\]
 \[c_2+a_2+c_1=a_1.\]Substituting $a_1 = c_1/e$ and $a_2 = c_2/e$, we get
 \[c_2 + \frac{c_2}{e} + c_1=\frac{c_1}{e}\]
 \[c_2e + c_2 + c_1e = c_1\]
 \[c_2(e+1) = c_1(1-e)\]
 \[\frac{c_1}{c_2} = \frac{1+e}{1-e}\]We now take into account the information regarding the ratio of the ellipses' areas. Because the area of an ellipse with semi-major axis $2a$ and eccentricity $e$ can be written as $\pi a^2\sqrt{1-e^2}$, we get that:
 \[\frac{\text{Area}_1}{{\text{Area}_2}} = \frac{a_1^2\sqrt{1-e^2}}{a_2^2\sqrt{1-e^2}} = \frac{a_1^2}{a_2^2} = 2025.\]

 Therefore, $\frac{a_1}{a_2} = 45$. Because $a_1 = c_1/e$ and $a_2 = c_2/e$, then $\frac{c_1/e}{c_2/e} = \frac{c_1}{c_2}=45$ as well.

 Substituting this into the equation found earlier, we get that 
 \[\frac{c_1}{c_2} = \frac{1+e}{1-e} = 45\]
 \[1+e = 45-45e\]
 \[46e = 44\]
 \[e = \boxed{\frac{22}{23}}.\]




\textbf{Solution 3 (No area formulas!)}

In the larger ellipse, let $a$ be half the major axis (also called the semi-major axis) and $b$ be half the minor axis (also called the semi-minor axis), and $c$ be the distance between its center and a focus. Since we are given that $2025 = 45^2$ is the ratio of the areas of the two ellipses, it follows that the linear scale factor of the two is $45:1$. Therefore, the semi-major axis of the smaller ellipse is $\frac{1}{45}a$ and its semi-minor axis is $\frac{1}{45}b$. Furthermore, the distance between a focus of the smaller ellipse and its center is $\frac{1}{45}c$. There are now two ways to express the length of the major axis of the large ellipse: $2a$ and $a + c + \frac{1}{45}c + \frac{1}{45}a$. Setting these equal, we now have an expression of degree 1 entirely in terms of $a$ and $c$. Some algebra gives us:

 

 \[a-\frac{1}{45}a=c+\frac{1}{45}c\]

 \[\frac{44}{45}a=\frac{46}{45}c\]

 \[44a=46c\]

 \[e=\frac{c}{a}=\boxed{\text{(D) }\dfrac{22}{23}}\].

 
Note that this method did not require you to know the area formula of an ellipse (which, incidentally, is $\pi a b$).

 



\textbf{Solution 4 (Similarity)}

Parametrize the bigger ellipse with $a,c$ as defined in the problem, and the smaller ellipse with $a',c'$. Note that the ellipses have the same eccentricity, so they are similar with ratio $k^2=2025$ or $k=45$. We wish to find $e=\frac{c}{a}=\frac{c'}{a'}$.

 Looking at the figure, we can decompose the semi-major axis $a$ of the bigger ellipse as
 \[a=c+c'+a'\]Using $c=kc'=45c'$ and $a=45a'$ and rearranging gives us:
 \[45a'-a'=44a'=45c'+c'=46c'\]Thus $e=\frac{c'}{a'}=\frac{44}{46}=\boxed{\textbf{(D)}~\frac{22}{23}}.$



	\item (\textit{AMC 12 12A 2025}) What is the value of

 \[\log_{2}\left( \sin\left( \frac{\pi}{8} \right)  \right) + \log_{2}\left( \sin\left( \frac{3\pi}{8} \right)  \right) + \log_{2}\left( \sin\left( \frac{5\pi}{8} \right)  \right) + \log_{2}\left( \sin\left( \frac{7\pi}{8} \right)  \right)?\]
$\textbf{(A)}\   -3  \qquad                 \textbf{(B)}\   -2\sqrt{2}  \qquad                 \textbf{(C)}\   -2  \qquad                 \textbf{(D)}\   -\frac{3}{2}  \qquad                 \textbf{(E)}\   -\sqrt{2}$

 \textbf{Answer: }C



\textbf{Solution 1}

Let our subset be $\{11,12,13,...,20\}.$ If we add any element from the set $\{1,2,3,...,10\}$ to our current subset, we will have to remove at least one element from our subset. We cannot add 10 because if we do, the sum of 10 and itself is 20. Hence, the maximum size of our subset is $\boxed{\text{(C) }10}$.

 



\textbf{Solution 2 (Quicksolve)}

Notice that an odd number plus an odd number always results in an even number. Consider the subset of all odd numbers $S = \{1, 3, 5, ..., 19\}$. The addition of any even number will result in a violation of the rules, so the maximum number of elements in this subset is $\boxed{\text{(C) }10}$.

 \textbf{Note (Those who know)} This problem is essentially the same as \href{https://artofproblemsolving.com/wiki/index.php?title=2022\_AMC\_10B\_Problems/Problem\_14}{2022 AMC 10B Problem 14}.

 



\textbf{Solution 3}

Let $S = \{1, 2, 3, ..., 20\}$. We are looking for the largest possible sum-free subset $A \subseteq S$.

 First, we can show that a sum-free set of size 10 is possible.Consider the set $A_1 = \{11, 12, ..., 20\}$. The smallest possible sum of two elements (not necessarily distinct) from this set is $11 + 11 = 22$. Since $22 > 20$, no sum of two elements from $A_1$ can be in $A_1$. The size of $A_1$ is $20 - 11 + 1 = 10$.Thus, a sum-free set of size 10 does exist.

 Next, we prove that no sum-free set $A \subseteq S$ can have 11 or more elements.Let $m$ be the largest element in $A$.Since $A$ is sum-free and $m \in A$, for any $x \in A$, we must have $m-x \notin A$ (unless $x = m-x$, in which case $x \notin A$). This is because if both $x$ and $m-x$ were in $A$, their sum $x + (m-x) = m$ would be in $A$, which is a contradiction.

 We can partition the set $\{1, 2, ..., m\}$ into pairs that sum to $m$, plus any leftover elements.

 Case 1: $m$ is even. Let $m = 2k$.The numbers in $\{1, ..., m\}$ can be grouped as:$\{ (1, 2k-1), (2, 2k-2), ..., (k-1, k+1) \}$ along with the numbers $k$ and $m=2k$.

 
            \begin{itemize}
        \item $A$ contains $m$.
        \item $A$ cannot contain $k$, because $k+k = 2k = m$, which is in $A$.
        \item $A$ can contain at most one number from each of the $k-1$ pairs.
            \end{itemize}
            The maximum number of elements $A$ can contain from $\{1, ..., m\}$ is $1 + 0 + (k-1) = k = m/2$.Since $m \le 20$, the maximum size of $A$ is $m/2 \le 20/2 = 10$.

 Case 2: $m$ is odd. Let $m = 2k-1$.The numbers in $\{1, ..., m\}$ can be grouped as:$\{ (1, 2k-2), (2, 2k-3), ..., (k-1, k) \}$ along with $m=2k-1$.

 
            \begin{itemize}
        \item $A$ contains $m$.
        \item $A$ can contain at most one number from each of the $k-1$ pairs.
            \end{itemize}
            The maximum number of elements $A$ can contain from $\{1, ..., m\}$ is $1 + (k-1) = k$.Since $m = 2k-1$, we have $k = (m+1)/2$.The largest possible odd $m$ is 19. If $m=19$, $k = (19+1)/2 = 10$.So the maximum size of $A$ is $(m+1)/2 \le (19+1)/2 = \boxed{\textbf{(C) } 10}$.

 In all cases, a sum-free subset of $S$ can have at most 10 elements.Since we found a set with 10 elements, the greatest possible number of elements is $10$.

 



\textbf{Solution 4 (Another quicksolve)}

Note that selecting integers from $\{11, 12, ..., 20\}$ also satisfy the constraints of the problem. The minimum sum of any two members of the set is $11 + 11 = 22$ which is greater than any of the elements.

 



\textbf{Solution 5 (with quick proof)}

As all of the above solutions point out, the set $\{11, ..., 20\}$ satisfies the required condition, and so the greatest size of a sum-free set is no less than 10. How do we know that 11 isn't possible? Induction is our friend here. 

 We want to prove that the set $\{1, ..., 2n\}$ has a sum-free set of size no more than $n$. It is obvious for $n=1$. Suppose it is true for $n-1$. Now let $S$ be a sum-free subset of $\{1, ..., 2n\}$ and that it has at least $n+1$ elements. The intersection of $S$ and $\{1, ..., 2n-2\}$ is sum-free, so it has no more than $n-1$ elements. So we must include both $2n-1$ and $2n$ into $S$. If so, then $1$ and $n$ cannot belong to $S$. 

 We are now searching for $n-1$ more elements. But we can make $n-2$ 'pigeonholes': $\{2, 2n-2\}, \{3, 2n-3\}, ..., \{n-1, n+1\}$, seeing that they can provide no more than $n-2$ elements because each pigeonhole will have a sum of $2n$.

 


 


	\item (\textit{AMC 12 12A 2025}) Annie, Bobby, and Casey each receive an integer, not necessarily distinct, from $1$ through $10$. In how many ways can the numbers be given so that their sum is divisible by $7$?


 $\textbf{(A)}\   139  \qquad                  \textbf{(B)}\   140  \qquad                 \textbf{(C)}\   142  \qquad                 \textbf{(D)}\   143  \qquad                 \textbf{(E)}\   145$

 \textbf{Answer: }D



\textbf{Solution 1}

Let $CD \perp AB$ with foot $D$. Right triangles $ACD$ and $BCD$ give $AC^2 = AD^2+CD^2$, $BC^2 = BD^2+CD^2$, $AC^2-BC^2 = AD^2-BD^2 =(AD-BD)(AD+BD)$. 

 Since $AD+BD = AB = 80$ and $AC^2-BC^2 = 75^2-45^2 = 3600$, we get the equation $3600 = 80(AD-BD)$. This equation simplifies to $45 = AD - BD$. We can solve the system of equations $AD + BD = 80$ and $AD - BD = 45$ easily via elimination, and we get $AD = \frac{125}{2}$, $BD = \frac{35}{2}$. $CD^2 = AC^2-AD^2 = 75^2-\left(\frac{125}{2}\right)^2 = \frac{6875}{4}$, $CD = \frac{25\sqrt{11}}{2}$. 

 By Angle Bisector Theorem, $\frac{DP}{PC} = \frac{DB}{BC} = \frac{\frac{35}{2}}{45} = \frac{7}{18}$, $PC = CD-DP$ thus, $18DP = 7(CD-DP)$, $25DP = 7CD$, $DP = \left(\frac{7}{25}\right)CD = \left(\frac{7}{25}\right)\left(\frac{25\sqrt{11}}{2}\right) = \frac{7\sqrt{11}}{2}$. $BP^2 = BD^2+DP^2 = \left(\frac{35}{2}\right)^2+\left(\frac{7\sqrt{11}}{2}\right)^2 = \frac{1225}{4}+\frac{49(11)}{4} = \frac{1764}{4} = 441$, thus $BP = \boxed{\text{(D) }21}.$




\textbf{Solution 2 (Law of Cosines)}

Scale this down to a $9-15-16$ triangle (we will multiply the result by $5$ in the end).Note that $9^2+16^2-18\cdot16\cdot\cos(\angle B) = 15^2$, so $112 = 18 \cdot 16 \cdot \cos(\angle B)$, which simplifies to $\frac{7}{18} = \cos(\angle B)$. Then $\cos(\frac{\angle B}{2}) = \sqrt{\frac{1+\frac{7}{18}}{2}} = \sqrt{\frac{25}{36}} = \frac{5}{6}$ (positive root since the angle is acute). Therefore, we have $\frac{BD}{BP} = \frac{5}{6}$, assuming that $D$ is the foot of the altitude.There are many ways to proceed from here to find $BD$. Note that by Heron's formula, the area of the scaled-down triangle is $\sqrt{(20)(4)(5)(11)} = 20\sqrt{11}$. Therefore, $CD = \frac{5}{2}\sqrt{11}$. Using Pythagorean Theorem, we get $BD= \sqrt{81-\frac{25 \cdot 11}{4}} = \sqrt{\frac{324-275}{4}} = \frac{7}{2}$. Therefore, we get $\frac{\frac{7}{2}}{BP} = \frac{5}{6}$, so $BP = \frac{21}{5}$, and we scale up by $5$ to get $\boxed{\text{(D) }21}$.~ScoutViolet

 



\textbf{Solution 3 (Stewarts)}

Let the foot of the altitude coming from $C$ on segment $AB$ be $D$. Using the fact that $CD$ is a common leg in right triangles $\triangle CDB$ and $\triangle CDA$, we have 
 \[45^2 - BD^2 = 75^2 - (80-BD)^2.\] Expanding gives 
 \[45^2=75^2-80^2+160BD,\] so $BD=\frac{35}{2}.$ Let the foot of the angle bisector from $B$ to $AC$ be point $E$. Since $BE$ is the angle bisector of $\angle CBA$, we can use the angle bisector theorem. This gives 
 \[\frac{45}{CE}=\frac{80}{75-CE},\] so $CE=27$ and $AE=75-27=48$. Now we can use Stewart’s Theorem to find $\overline{BE}$. We have 
 \[(48)(45)^2+(27)(80)^2=(27)(48)(75)+75BE^2.\] To simplify this expression, just divide by the greatest common divisor and solve from there. In the end, we get $BE=48$. Let $BP=x$, so $PE=48-x$. Draw the altitude from $E$ down to $AB$. Let the foot of this altitude be $F$. Since $EF || CD$, we have $\triangle AFE \sim \triangle ADC$. Hence, we can write the equation 
 \[\frac{48}{75}=\frac{AF}{\frac{125}{2}}.\] Solving gives $AF=40$, so $FD=\frac{125}{2}-40=\frac{45}{2}$. Since $PD || EF$, we also have $\triangle BDP \sim \triangle BFE$, so we have 
 \[\frac{x}{48}=\frac{\frac{35}{2}}{40}.\] Solving for $x$ gives $\boxed{x=21}$ or $\boxed{\text{D}}.$




\textbf{Solution 4 (Rulerbash)}

Preface: When we have a problem as such, involving a simple diagram with minimal instructions, I use a method I named "rulerbash". Rulerbash should only be used in specific cases and as a last resort, mainly in the event of a time crunch or a difficult problem. 

 Start with the longest side, drawing a line with a length of $8\,\text{cm}$ ($AB$). Then, using a compass, draw 2 circles centered around points $A$ and $B$, $7.5\,\text{cm}$ and $4.5\,\text{cm}$ radii respectfully. At the point of intersection of these 2 circles, we have point $C$, completing a perfectly scaled drawing of $\triangle ABC$. (Note the circles are not necessary with a bit of trial and error with the side lengths, they simply offer a way to get it done first try). 

 
                    \begin{figure}[H]
                        \centering
                        \includegraphics[width=0.5\textwidth]{Solution/AMC_12_12A_2025_16_Solution_Figure_1.png}
                    \end{figure}
                    

 


 
After this, dropping the altitude to $AB$ is simple with a ruler and careful placement, and angle bisector can be estimated quite accurately. (Note: Alternately, a compass can be used to accurately place the altitude and the bisector, which are both trivial constructions ~ZingiberiMaracandae)

 
                    \begin{figure}[H]
                        \centering
                        \includegraphics[width=0.5\textwidth]{Solution/AMC_12_12A_2025_16_Solution_Figure_2.png}
                    \end{figure}
                    

 After all of this, we can reuse our ruler and measure $BP = 2.1\,\text{cm}$, and using our scale of $80=8\,\text{cm}$, our final answer is $\boxed{\text{(D) }21}.$

 Other note: To learn how to draw a angle bisector of an angle to to \href{https://www.mathsisfun.com/geometry/construct-anglebisect.html}{https://www.mathsisfun.com/geometry/construct-anglebisect.html}

 



\textbf{Solution 5 (No Trig)}

Let \(D\) be the intersection of the altitude from \(C\) to \(AB\). To simplify calculations, divide all side lengths by \(5\), and multiply by \(5\) again at the end.  First, we use Heron’s Formula, \(\sqrt{s(s-a)(s-b)(s-c)}\), to find the area. Let \([ABC]\) denote the area of \(\triangle ABC\). By Heron’s Formula,  
 \[[ABC] = \sqrt{20 \cdot 5 \cdot 4 \cdot 11} = 20\sqrt{11}.\]  Next, we find the altitude \(CD\) using the formula for the area of a triangle, \(\frac{1}{2}bh = \text{area}\):  
 \[\frac{1}{2}(16)(CD) = 20\sqrt{11} \quad \Rightarrow \quad CD = \frac{5\sqrt{11}}{2}.\]  We can use the Pythagorean Theorem in \(\triangle CDB\) to find \(DB\):  
 \[DB^2 + \frac{25 \cdot 11}{4} = 81 \quad \Rightarrow \quad 4DB^2 + 275 = 324 \quad \Rightarrow \quad DB^2 = \frac{49}{4},\]so \(DB = \frac{7}{2}\).  Next, we use the Angle Bisector Theorem to find \(PD\). Let \(x = PC\) and \(y = PD\). Since \(x + y = \frac{5\sqrt{11}}{2}\), we have \(x = \frac{5\sqrt{11}}{2} - y\).  From the given ratio,  
 \[\frac{9}{x} = \frac{7}{2y} \quad \Rightarrow \quad 18y = 7x.\]  Substituting \(x = \frac{5\sqrt{11}}{2} - y\),  
 \[18y = 7\left(\frac{5\sqrt{11}}{2} - y\right) \quad \Rightarrow \quad 25y = \frac{35\sqrt{11}}{2},\]so \(y = \frac{7\sqrt{11}}{10}\).  

 Now, using the Pythagorean Theorem again to find \(BP\):  
 \[\frac{49 \cdot 11}{100} + \frac{49}{4} = BP^2 \quad \Rightarrow \quad 100BP^2 = 49(11 + 25) = 49 \cdot 36,\]so \(BP = \frac{42}{10} = \frac{21}{5}.\)  

 Finally, multiplying the side lengths by \(5\) again gives \(BP = 21.\), or $\boxed{\text{D}}.$




\textbf{Solution 6 (Similar Triangle)}


                    \begin{figure}[H]
                        \centering
                        \includegraphics[width=0.5\textwidth]{Solution/AMC_12_12A_2025_16_Solution_Figure_3.png}
                    \end{figure}
                    

 Due to Angle Bisector Theorem
 \[\frac{CE}{AE} = \frac{BC}{BA} = \frac{45}{80} = \frac{9}{16}\]
 \[CE = AC \times \frac{9}{25} = 27\]
 \[AE = AC \times \frac{16}{25} = 48\]Notice that $\dfrac{CE}{CB} = \dfrac{CB}{CA} = \dfrac{3}{5}$, $\triangle CBE \sim \triangle CAB$, $\angle A = \angle CBE = \angle EBF$, so $\triangle ABE$ is isosceles triangle, hence
 \[AF = BF\]Notice that $AF:FD = AE:EC = 16:9$, $DB = FB - FD$, then 
 \[AF:DB = 16:7\]

 Notice that $\triangle PBD \sim \triangle EAF$, $PB = AE \times \dfrac{BD}{AF} = 48 \times \dfrac{7}{16} = \boxed{\text{(D) }21}.$




\textbf{Solution 7 (Pythagoras+angle bisector theorem)}

Let $CD$ be the altitude to side $AB$, and let $E$ be the point of intersection of the angle bisector of $\angle B$. Let $P$ be the point of intersection of the angle bisector and the altitude.

 By Pythagoras, $CD^2=BC^2-BD^2=AC^2-AD^2$

 We let $BD = x$, so $AD=80-x$

 Now we have

 
 \[45^2-x^2=75^2-(80-x)^2\]

 (This looks complicated, but we can see that we are going to have two $-x^2$ terms on both sides, so it cancels out)

 Simplifying and rearranging gives us $80-2x=45$, so $x=\frac{35}{2}$

 Therefore, $BD=\frac{35}{2}$, and $AD=\frac{125}{2}$. This gives us $CD=\frac{25}{2}\sqrt{11}$

 
Now, note that in triangle $BDC$, $BP$ bisects $\angle DBC$

 By the angle bisection theorem, we have$\frac{BD}{DP}=\frac{BC}{CP}$

 which gives us $\frac{DP}{CP}=\frac{BD}{BC}=\frac{7}{18}$

 We have $DP=\frac{7}{25}*CD=\frac{7}{2}\sqrt{11}$ and $CP=\frac{18}{25}*CD=\frac{18}{2}\sqrt{11}$

 So

 $BP^2=BD*BC-DP*CP=\frac{35}{2}*45-\frac{7}{2}*\frac{18}{2}*11=441$

 Therefore, $BP=21$

 (This is my first solution, pls point out any mistakes in math or LATEX). ~backtosq-1

 \textbf{Alternative Ending to Solution 7} As an alternative ending equation to the solution above, I used the Pythagorean Theorem again, giving us: $BP^2=BD^2+PD^2=(\dfrac{35}{2})^2+(\dfrac{7\sqrt{11}}{2})^2=(\dfrac{1225}{4})+(\dfrac{539}{4})=\dfrac{1764}{4}=441$

 Thus, $BP=\boxed{\textbf{(D)}~21}$




\textbf{Solution 8 (Pythagora's + Double Angle Identity)}

Since all the lengths share a factor of $5$, divide by $5$ and multiply it back at the end for smaller numbers. Drop the altitude from vertex $C$ to point $D$ on $AB$. Then comparing the altitude by Pythagora's on both sides gives 
 \[9^2-x^2=15^2-(16-x)^2=15^2-16^2+32x-x^2=-31+32x-x^2\]Cancelling the $-x^2$ then solving gives us $x=\frac{112}{32}=\frac{7}{2}$. Since $\triangle BDC$ is a right triangle with $\angle B=2\theta$, we find $\cos(2\theta)=\frac{\frac{7}{2}}{9}=\frac{7}{18}$.

 Let $\theta=\angle PBA=\angle PBC$. Because $\angle B$ is opposite to the leg $9<15$ and not the hypotenuse, we must have that $2\theta$ is acute, likewise for $\theta$. Hence, $\cos \theta>0$.

 Thus, using cosine definition on $\triangle BDP$ and the double angle identity $\cos 2\theta=2\cos^2 \theta-1$, we get
 \[\cos \theta=\frac{7/2}{BP}=\sqrt{ \frac{1+\cos 2\theta}{2} }=\sqrt{ \frac{1+\frac{7}{18}}{2} }=\sqrt{ \frac{25}{36} }=\frac{5}{6}\]Finally, solving for $BP=\frac{7}{2\cos\theta}=\frac{7}{2}\cdot \frac{6}{5}=\frac{21}{5}$. Scaling back to the original diagram by $5$ gives $\boxed{\textbf{(D)}~21}$.

 



\textbf{Solution 9(Law of Cosines+Angle Bisector Length Theorem)}

Use Law of Cosines in $\triangle ABC$:
 \[\cos B=\frac{80^2+45^2-75^2}{2\cdot 80\cdot 45}=\frac{7}{18}\]Call $E$ the intersection of the angle bisector of $B$ and line segment $AC$, and let $CD \perp AB$, $D=CD\cap AB$, therefore $P=CD\cap BE$. In right triangle $\triangle BCD$, use the definition of cosine:
 \[\cos B=\frac{BD}{BC}=\frac{7}{18},\sin B=\sqrt{1-\cos^{2}B}=\frac{5\sqrt{11}}{18}\]Therefore 
 \[BD=\frac{35}{2}, CD=\frac{25\sqrt{11}}{2}\]Use Angle Bisector Theorem in $\triangle BCD$: 
 \[\frac{BD}{BC}=\frac{PD}{PC}=\frac{7}{18}\]Therefore, 
 \[PD=\frac{7}{25}CF=\frac{7\sqrt{11}}{2}, PC=\frac{18}{25}CF=9\sqrt{11}\]Use Angle Bisector length formula:
 \[BP=\sqrt{BD\cdot BC-PD\cdot PC} =\boxed{\textbf{(D)} \ 21}\]




\textbf{Solution 10}

Let the point where the altitude to $AB$ meets $AB$ be $D$. We can compute $D$ easily. Let $BD = a \implies AD = 80 - a$. So

 $45^2 - a^2 = 75^2 - (80 - a)^2 \implies 45^2 - a^2 = 75^2 - 80^2 + 160a - a^2 \implies 45^2 - 75^2 + 80^2 = 160a \implies -30 \cdot 120 + 80 \cdot 80 = 10 \cdot 16a \implies -360 + 640 = 16a \implies a = \frac{280}{16} = \frac{35}{2}$. 

 Now let $\angle DBP = \theta$. Note this is an angle bisector so we can find $\cos(\angle ABC)$ where $\theta = \frac{1}{2} \angle ABC$. So we use Law of Cosines to get

 $80^2 + 45^2 - 160 \cdot 45 \cos(2\theta) = 75^2 \implies 16^2 + 9^2 - 32 \cdot 9 \cdot \cos(2\theta) = 15^2 \implies \cos(2\theta) = \frac{112}{32 \cdot 9} = \frac{7}{18}$. Now note

 $\cos(2\theta) = 2\cos^2(\theta) - 1 = \frac{7}{18} \implies \cos^2(\theta) = \frac{25}{36} \implies \cos(\theta) = \frac{5}{6}$. 

 Finally, by $SOH-CAH-TOA$ note that $\cos(\theta) = \frac{BD}{BP} = \frac{5}{6}$ and we know $BD = \frac{35}{2}$ so 

 $BP = \frac{6}{5} \cdot \frac{35}{2} = \boxed{21}$.

 


 


	\item (\textit{AMC 12 12A 2025}) Belinda follows a path on the plane that is entirely contained within the square region $0 \leq x, y \leq 4$. Call her path \textit{dignified} if and only if all of the following conditions are satisfied.


 
            \begin{itemize}
        \item Each step is one unit up, down, or to the right.

        \item The path does not revisit any vertex.

        \item The path starts at $(a, 0)$ and ends at $(b,0)$, where a and $b$ are integers with $a < b$, and touches the $x$-axis only at these two points.
            \end{itemize}
            The following figure shows some examples of dignified and not dignified paths.


 
                    \begin{figure}[H]
                        \centering
                        \includegraphics[width=0.5\textwidth]{Problem/AMC_12_12A_2025_17_Problem_Figure_1.png}
                    \end{figure}
                    


 How many dignified paths can Belinda follow?


 $\textbf{(A)}\   174  \qquad                 \textbf{(B)}\   340  \qquad                 \textbf{(C)}\   448  \qquad                 \textbf{(D)}\   512  \qquad                 \textbf{(E)}\   970$

 \textbf{Answer: }A



\textbf{Solution 1 (substitution)}

We initially note the symmetry of the product $(z+i)(z+2i)(z+3i)$, in that the factors $(z+i)$ and $(z+3i)$ can be expressed as $(z+2i) \mp i$ respectively. Accordingly, making the substitution $w = z+2i$ enables us to apply the 'difference of two squares' identity, yielding 
 \begin{align*}(z+i)(z+2i)(z+3i)+10 &= (w-i)(w)(w+i)+10 \\ &= w\left(w^2-i^2\right)+10 \\ &= w\left(w^2+1\right)+10 = w^3+w+10.\end{align*}
Having thus eliminated the quadratic term of the given cubic polynomial, we easily observe that $-2$ is a root of this new cubic (in particular, by the \href{https://artofproblemsolving.com/wiki/index.php?title=Rational\_root\_theorem}{rational root theorem}, the only possible rational roots are integers that are factors of $10$).

 It follows that $w^3+w+10 = \left(w+2\right)\left(w^2-2w+5\right)$, and hence, using either the quadratic formula or completing the square, the other two roots are $1+2i$ and $1-2i$. Moreover, since our initial substitution simply shifts all points in the complex plane by $2i$ (i.e. by $2$ units vertically upward) and area is preserved under translation, the required area of the triangle in the $z$-plane is the same as that of the triangle formed by these three roots in the $w$-plane.

 Now, as $1+2i$ and $1-2i$ of course have the same real part, the side of the triangle between these two vertices is simply vertical, and has length $2-(-2) = 4$. The length of the altitude to that side is the perpendicular distance from $-2$ to the vertical line $\usepackage{physics}\Re(z) = 1$, i.e. the horizontal distance from $-2$ to $1$, which is $1-(-2) = 3$. Thus the area of the triangle is $\frac{1}{2} \cdot 4 \cdot 3 = \boxed{\textbf{(A)}~6}$.

 



\textbf{Solution 1.1 (alternative using shoelace formula for last step)}

While the area of the triangle (once we have determined its vertices) can be straightforwardly computed using $\text{area } = \frac{1}{2} \cdot \text{base} \cdot \text{perpendicular height}$ as explained above, it is also possible to use the \href{https://artofproblemsolving.com/wiki/index.php?title=Shoelace\_Theorem}{shoelace formula}, a general result giving the area of any polygon in the plane in terms of the coordinates of its vertices.

 Specifically, the vertices $-2$, $1+2i$, and $1-2i$ in the complex plane of course become the points $(-2,0)$, $(1,2)$, and $(1,-2)$ in the $x$-$y$ plane, so the shoelace formula gives the area of the triangle as 
 \[\frac{1}{2} \cdot \left\lvert (-2) \cdot 2 + 1 \cdot (-2) + 1 \cdot 0 - 0 \cdot 1 - 2 \cdot 1 - (-2) \cdot (-2)\right\rvert = \frac{1}{2} \cdot \left\lvert -12\right\rvert = \boxed{\textbf{(A)}~6}.\]




\textbf{Solution 1.2 (alternative using triangle in -plane for last two steps)}

As an alternative to arguing that the triangle formed by the three roots for $w$ has the same area as the one originally formed by the roots for $z$, we can simply convert back to the $z$-plane using $z = w-2i$. Therefore, as the roots for $w$ are $-2$, $1+2i$, and $1-2i$, those for $z$ are $-2-2i$, $(1+2i)-2i = 1$, and $(1-2i)-2i = 1-4i$, with the vertical side of the triangle thus being that between $1$ and $1-4i$.

 We can now use either of the two approaches above: for the former, the vertical side again has length $0-(-4) = 4$ and the corresponding altitude has length $1-(-2) = 3$, once again giving the area as $\frac{1}{2} \cdot 4 \cdot 3 = \boxed{\textbf{(A)}~6}$. For the latter, the shoelace formula also gives the area as 
 \[\frac{1}{2} \cdot \left\lvert (-2) \cdot 0 + 1 \cdot (-4) + 1 \cdot (-2) - (-2) \cdot 1 - 0 \cdot 1 - (-4) \cdot (-2)\right\rvert = \frac{1}{2} \cdot \left\lvert -12\right\rvert = \boxed{\textbf{(A)}~6}.\]




\textbf{Solution 2 (directly finding the roots)}

Expanding the given polynomial yields 
 \begin{align*}(z+i)(z+2i)(z+3i)+10 &= \left(z^2 + 2iz + iz + 2 \cdot (-1)\right)\left(z+3i\right)+10 \\ &= \left(z^2+3iz-2\right)\left(z+3i\right)+10 \\ &= z^3 + 3iz^2 + 3iz^2 + 9 \cdot (-1)z - 2z - 6i + 10 \\ &= \left(z^3-11z+10\right)+\left(6z^2-6\right)i,\end{align*}
 and so any complex number $z$ will be a root if and only if it is a root of both the real part, $\left(z^3-11z+10\right)$, and the imaginary part, $\left(6z^2-6\right)$.

 We immediately observe that $1$ is indeed a root of both of them, so by factoring out $(z-1)$, we obtain $(z+i)(z+2i)(z+3i)+10 = \left(z-1\right)\left(z^2+(1+6i)z+(-10+6i)\right)$, and now using the quadratic formula, we deduce that the other two roots are given by 
 \begin{align*}z &= \frac{-(1+6i) \pm \sqrt{(1+6i)^2 - 4 \cdot 1 \cdot (-10+6i)}}{2 \cdot 1} \\ &= \frac{-1-6i \pm \sqrt{1+12i-36+40-24i}}{2} \\ &= \frac{-1-6i \pm \sqrt{5-12i}}{2}.\end{align*}


 Noting that $5-12i = (3-2i)^2$ (or else by writing $5-12i = (a+bi)^2$ for some $a,b \in \mathbb{R}$, and then solving for $a$ and $b$ by comparing real and imaginary parts of both sides), the above expression reduces to $z = \frac{-1-6i+3-2i}{2} = \frac{2-8i}{2} = 1-4i$ or $z = \frac{-1-6i-3+2i}{2} = \frac{-4-4i}{2} = -2-2i$. In other words, the three roots are precisely $1$, $1-4i$, and $-2-2i$, and thus using either of the approaches in Solution 1.2 above, the area of the triangle is $\boxed{\textbf{(A)}~6}$.

 


	\item (\textit{AMC 12 12A 2025}) A ribbon wraps around a pole in such a way that equal amounts of area are covered and not covered by the ribbon. The width of the ribbon is $4$ cm, and the circumference of the pole is $10$ cm. The figure below shows a possible front view of a portion of the pole.


 
                    \begin{figure}[H]
                        \centering
                        \includegraphics[width=0.5\textwidth]{Problem/AMC_12_12A_2025_18_Problem_Figure_1.png}
                    \end{figure}
                    


 In centimeters, how much of the pole height is needed for the ribbon to wrap around exactly once?


 $\textbf{(A)}\   \frac{15}{2}  \qquad                 \textbf{(B)}\   8  \qquad                 \textbf{(C)}\   10  \qquad                 \textbf{(D)}\   \frac{25}{2}  \qquad                 \textbf{(E)}\   \frac{40}{3}$

 \textbf{Answer: }C



\textbf{Solution 1}

Let $0<x<y<z \le 8$;$x$ cannot be $1$ because it makes $xy>z$ $\rightarrow$ $y>z$;

 $x=2$, $y=3$, $z$ can be $4$, $5$ but not others; $x=2$, $y=4$, $z$ can be $5$, $6$, $7$;$x=2$, $y=5$, $z$ can be $6$, $7$, $8$;$x=2$, $y=6$, $z$ can be $7$, $8$;$x=2$, $y=7$, $z$ can be $8$;for $x=2$, total $11$ cases; 

 Similarly, for $x=3$, $y=4$, $5$, $6$, $7$, total $10$ cases; for $x=4$, $y=5$, $6$, $7$, total $6$ cases; $x=5$, $y=6$, $7$, $3$ cases; $x=6$, $y=7$, $z=8$, $1$ cases;

 Total $= 11 + 10 + 6 + 3 + 1 = 31$. Permutate $x$, $y$, $z$ for ordered triple, it is $31 \cdot 6=186$, $\fbox{C}$.

 



\textbf{Solution 2}

For now, assume $x\leq y\leq z$.

 First note that no number can be 0, as it would imply $0y>z$. Similarly, no number can be 1, as it would imply $1y>z$. So we only need to consider numbers between 2 and 8, inclusive.

 We may use complementary counting:

 Consider when $xy\leq z$. This implies $xy\leq 8$. Some quick calculations gives us the products $2\cdot2,2\cdot3,2\cdot4$. We may now calculate the number of times each happens (we are no longer assuming $x\leq y\leq z$):

 
            \begin{itemize}
        \item $2\cdot2$: This case is invalid as it asks for distinct integers.
        \item $2\cdot3$: $2,3,(n\geq6)$. Then we have $6\cdot3=18$ cases.
        \item $2\cdot4$: $2,4,(n\geq8)$. Then we have $6\cdot1=6$ cases.
            \end{itemize}
            In total, there are $7\cdot6\cdot5=210$ total cases, so our final answer is $210 - (18 + 6) = \boxed{186}$.

 



\textbf{Solution 3}

Note that if $x,y,z \ge 3$, all pairs work - hence, we have $\binom{6}{3} \cdot 6 = 120$ pairs.Now, note that if $x=1$, we get $y>z$ and $z>y$, contradiction. Therefore, assume $x=2$ (since we said $x < 3, x \neq 1$). Note that we need $2y>z>\frac{y}{2}$. WLOG assume $y>z$, we get there are $11$ pairs that work (we just need $2z>y$): $(y,z) = (8,5), (8,6), (8,7), (7,4), (7,5), (7,6), (6,4), (6,5), (5,3), (5,4), (4,3)$. With $x=2$, we can re-arrange these in $3!=6$ ways each, hence the answer is just $120 + 6 \cdot 11 = \boxed{186}$.

 



\textbf{Solution 4 (Bounding)}

Note that if $x,y,z>2,$ then this must hold, as if it didn't, we would need, WLOG, $x>3y,$ which would imply $x>8.$ Therefore, there are at least $6\times 5\times 4=120$ solutions. However, we must have $x,y,z\ge2,$ as if any variable is $0,$ this clearly doesn't work, and $x=1$ gives $y>z$ and $z>y,$ which is impossible. Therefore, there are at most $7\times6\times5=210$ solutions. The only choice that works from here is $\boxed{\textbf{(C)}~186}$~krithikrokcs

 



\textbf{Solution 5 (Compl. Counting)}

The total number of cases is $7*6*5=210$. Note that no number can be 1, since that would result in cases like $z>y$ and $z<y$, which is impossible.

 We can try complementary counting. Since in this problem x,y,z cases are symmetrical, we first consider cases where $xy<z$. The cases are:$(x,y,z): (2,3,6) (2,3,7) (2,3,8) (2,4,8)$. By symmetry, the cases that don't work amount to $4*6=24$ cases. Therefore the total number of cases that work is $210-24=186$, which is $\boxed{\textbf{(C)}~186}$.~backtosq-1

 



\textbf{Solution 6 (Symmetry + Complementary Counting)}

Since $\{ x,y,z \} \subset \{ 1,2,3,4,5,6,7,8 \}$, WLOG let $1\le x<y<z\le 8$ ($\times 3! =6$). Note that the lower bounds make the conditions $y<xz\le 1z$ and $x<zy<z$ trivial. In addition, $x=1$ always fails, since it would force $z<xy=y$, a contradiction. Thus, it remains exactly to check $2\le x<y<z\le 8<xy$. 

 The only failure cases with $xy\le 8$ can occur when $(x,y)=(2,3),(2,4)$, for which we have corresponding $z\in \{ 6,7,8 \},\{ 8 \}$. The total number of cases is $\binom{7}{3}=35$. Thus, the number of good cases is $35-4=31$, and after multiplying by order we get $3!\cdot 31=\boxed{\textbf{(C)}~186}$.

 


	\item (\textit{AMC 12 12A 2025}) Starting with a list of three numbers, form a sequence in which each subsequent term is the average of the previous three terms. For example, starting with the sequence $(2,4,9)$, the next two terms are $\frac{2+4+9}{3} = 5$ and $\frac{4+9+5}{3} = 6$. Which of the following initial three terms will produce a sequence that has the greatest 2026th term?


 $\textbf{(A)}\   (1,2,3)  \qquad                 \textbf{(B)}\   (2,1,1)  \qquad                 \textbf{(C)}\   (3,4,8)  \qquad                 \textbf{(D)}\   (4,4,7)  \qquad                 \textbf{(E)}\   (4,5,6)$

 \textbf{Answer: }E



\textbf{Solution 1}

We begin by factoring:

 \begin{equation}\begin{split}&a^3b^2+a^2b^3+b^3c^2+b^2c^3+c^3a^2+c^2a^3 \\&=a^2b^2(a+b)+b^2c^2(b+c)+c^2a^2(c+a)\\&=(a^2b^2+b^2c^2+c^2a^2)(a+b+c) - a^2b^2c-b^2c^2a-c^2a^2b\\&=(a^2b^2+b^2c^2+c^2a^2)(a+b+c) - (abc)(ab+bc+ca).\\\end{split}\end{equation}From Vieta's Formulas, we know that $a+b+c = 0$, $ab+bc+ca = k$, and $abc = -1$. Therefore, the answer equals $(a^2b^2+b^2c^2+c^2a^2)(0) - (-1)(k) = \boxed{k}.$




\textbf{Solution 2}

Notice that this is $\sum_{\text{cyc}} a^2b^2(a+b)$.By Vieta's (just like in solution 1), we have $a+b+c=0$, so $a+b=-c$, $abc=-1$, and $ab+bc+ac=k$, so this is equivalent to $\sum_{\text{cyc}} -a^2b^2c = \sum_{\text{cyc}} -(abc)(ab) = \sum_{\text{cyc}} ab = k$ by Vieta's.~ScoutViolet

 



\textbf{Solution 3 (BS)}

Let’s just pick a nice and convenient a, b, c. Observe that $2\cdot2^{-\frac13}-2^\frac23=0$, so we can just have $(a,b,c)=(2^{-\frac13}, 2^{-\frac13},-2^\frac23$). So expanding $(x-a)(x-b)(x-c)$ we get $x^3-3\cdot2^{-\frac23}+1$. Evaluating $a^3b^2 + a^2b^3 + b^3c^2 + b^2c^3 + c^3a^2 + c^2a^3$, or $2a^5+2a^2c^3+2a^3c^2$, we get $-3\cdot2^{-\frac23}=k$, so the answer is $\boxed{\textit{E}}$. (i did this on the actual test)

 What’s up Ben ~ harrisontungtungtungsahur

 



\textbf{Solution 4 (Vieta's)}

First, note that $a+b+c=0$, $\sum_{cyc} ab = k$, $abc=-1$. We can factor the original expression into:

 
 \[\sum_{cyc} a^2b^2(a+b) = \sum_{cyc} \left(-\frac{1}{c}\right)^2(-c) = -\frac{\sum_{cyc} ab}{abc} = \boxed{\text{(E) } k.}\]




\textbf{Solution 5 (Newton's Sums) (If you suck at factoring like me)}

Note that $a^3b^2+a^2b^3+b^3c^2+b^2c^3+c^3a^2+c^2a^3 = (a^3+b^3+c^3)(a^2+b^2+c^2)-(a^5+b^5+c^5)$

 By Newton's sums, we have:
 \[(a^3 + b^3 + c^3) + k(a+b+c) + 3 = 0\]
 \[(a^5 + b^5 + c^5) + k(a^3+b^3+c^3) + 1(a^2+b^2+c^2) = 0\]

 Note that $a^2+b^2+c^2 = (a+b+c)^2 - 2(ab+bc+ca)$

 From Vieta's formulas, we have $a+b+c = 0$ and $ab + bc + ca = k$

 Solving, we get:
 \[(a^3 + b^3 + c^3) = -3\]
 \[(a^2 + b^2 + c^2) = -2k\]
 \[(a^5 + b^5 + c^5) - 3k - 2k = 0 \to (a^5 + b^5 + c^5) = 5k\]

 Therefore, $a^3b^2+a^2b^3+b^3c^2+b^2c^3+c^3a^2+c^2a^3 = (-3)(-2k)-5k = \boxed{\textbf{(E) }k}$




\textbf{Solution 6}

Like in solution 1, we begin by factoring:

 \begin{equation}\begin{split}&a^3b^2+a^2b^3+b^3c^2+b^2c^3+c^3a^2+c^2a^3 \\&=a^2b^2(a+b)+b^2c^2(b+c)+c^2a^2(c+a)\\\end{split}\end{equation}From Vieta's Formulas, we know that $a+b+c = 0$, $ab+bc+ca = k$, and $abc = -1$. 

 square $a+b+c = 0$ to get $a^2 + b^2 + c^2 = 0$, then isolating one variable each time yields:

 \begin{equation}\begin{split}&=-\frac{1}{c^2}(a+b)-\frac{1}{a^2}(b+c)-\frac{1}{b^2}(c+a)\\\end{split}\end{equation}

 Isolating a + b + c = 0, you get

 \begin{equation}\begin{split}&=-\frac{1}{c^2}(-c)-\frac{1}{a^2}(-a)-\frac{1}{b^2}(-b)\\&=-(\frac{ab+bc+ac}{abc})\\\end{split}\end{equation}

 This simplifies to:

 \begin{equation}\begin{split}&=-(\frac{k}{-1})\\&=k\end{split}\end{equation}

 Therefore, the answer equals = $\boxed{k}.$




\textbf{Solution 7 (Vieta + Factorization)}

We can easily factor the expression into 
 \[a^2b^2(a+b)+b^2c^2(b+c)+a^2c^2(a+c)\]

 By Vieta, we get $a+b+c=0$, so we can rewrite the expression into $a^2b^2(-c)+b^2c^2(-a)+a^2c^2(-b)$ which is just $-abc(ab+bc+ac)$. By Vieta we also get $abc=-1$ and $ab+bc+ac=k$. Therefore the expression is just $-(-1)k=$ $\boxed{k}.$




\textbf{Solution 8 (no Vieta + not quite factoring)}

Let the expression we are asked to evaluate equal $y$. 

 We have that $x^3+kx+1=(x-a)(x-b)(x-c)$. Expanding the right hand side yields $x^3-x^2(a+b+c)+x(ac+bc+ab)-abc$. Equating coefficients, $a+b+c=0$, $ac+bc+ab=k$, and $abc=-1$. 

 We may find the first equation helpful for cancelling things out, so let us fiddle around with it, trying to piece out as much of $y$ as possible: we discover that $(a+b+c)(ab)^2=a^3b^2+a^2b^3+c(ab)^2=0$, which shares terms with $y$. Likewise, $(a+b+c)(bc)^2=b^3c^2+b^2c^3+a(bc)^2=0$ and $(a+b+c)(ac)^2=a^3c^2+a^2c^3+b(ac)^2=0$. 

 Adding the three equations gives $y+a(bc)^2+b(ac)^2+c(ab)^2=0.$ Observe that $a(bc)^2+b(ac)^2+c(ab)^2=(abc)(ac+bc+ab)=-1(k)$. Thus $y-k=0$; $y=\boxed{k}$. 

 


 


	\item (\textit{AMC 12 12A 2025}) Which of the following statements most accurately describes the graph of the relationship

 \[\sin\left( x^{2} \right) = \cos\left( y^{2} \right)\]
$\textbf{(A)}\   \text{the set of all points } \left( \sqrt{ \frac{m}{\pi}}, \sqrt{\frac{n + \frac{1}{2}}{\pi}} \right), \text{ where } m \text{ and } n \text{ are nonnegative integers}$
$\textbf{(B)}\   \text{the set of all lines } y = \pm x + \frac{\pi}{2}, \text{ where } k \text{ is an integer}$
$\textbf{(C)}\   \text{the set of all circles } x^{2} + y^{2} = \frac{\pi}{2} + 2 \pi n, \text{ where n is a nonnegative integer}$
$\textbf{(D)}\   \text{a set of infinitely many concentric circles and infinitely many hyperbolas}$
$\textbf{(E)}\   \text{a set of finitely many ellipses centered at the origin and infinitely many parabolas}$

 \textbf{Answer: }C



\textbf{Solution 1 (Split Into Three Parts)}

Notice that the triangular faces have a slant height of $\sqrt{13^2-4^2}=\sqrt{153}$ and that the height is therefore $\sqrt{153-(\frac{13-7}{2})^2} = 12$. Then we can split the pentahedron into a triangular prism and two pyramids. The pyramids each have a volume of $\frac{1}{3}(3)(8)(12) = 96$ and the prism has a volume of $\frac{1}{2}(8)(12)(7) = 336$. Thus the answer is $336+96 \cdot 2 = \boxed{\textbf{(C) } 528}$




\textbf{Solution 2}

Note that the height is $12$ from the previous method.Note that as you go up, the length and width both decrease linearly and reach $0$ at the end.

 So the answer is $\int_0^{12} (8 - \frac{2}{3}x)(13 - \frac{1}{2}x) dx = \boxed{\textbf{(C) } 528}$

 -KapilTheAngel

 



\textbf{Solution (3D Pythagoras)}

Let $H$ be the height of the solid. By 3D Pythagoras (recursion of 2D Pythagoras),
 \[\left( \frac{8}{2} \right)^2+\left( \frac{13-7}{2} \right)^2+H^2=13^2=4^2+3^2+H^2=5^2+H^2\]This gives the $5-12-13$ triple, so $H=12$. Continue as in other solutions.

 \textbf{Note} It just reminds about one of old aime problem {\href{https://artofproblemsolving.com/wiki/index.php?title=1983\_AIME\_Problems/Problem\_11}{https://artofproblemsolving.com/wiki/index.php?title=1983\_AIME\_Problems/Problem\_11} 1983 AIME Problem 11}.~DRA777

 


	\item (\textit{AMC 12 12A 2025}) Let $T_{n}$ be the $n$th triangular number, that is, $T_{1} = 1$, $T_{2} = 1+2 = 3$, $T_{3} = 1+2+3 = 6$, and so on. The value of

 \[\sum_{k=1}^{\infty} T_{k} \cdot \left( \frac{1}{3} \right)^{k}\]
can be expressed as $\frac{p}{q}$, where $p$ and $q$ are relatively prime positive integers. What is $p+q$?


 $\textbf{(A)}\   14  \qquad                 \textbf{(B)}\   15  \qquad                 \textbf{(C)}\   16  \qquad                 \textbf{(D)}\   17  \qquad                 \textbf{(E)}\   18$

 \textbf{Answer: }A



\textbf{Solution 1}

The numerator can be written as $2^{2a} +2^{2a+2k}+...+2^{2a+2mk}$, which is actually a sum of geometric series. This can be expressed as $2^{2a} \cdot \frac{1-2^{2k(m+1)}}{1-2^{2k}}$. The denominator in the same way can be expressed as $2^a \cdot \frac{1-2^{k(m+1)}}{1-2^k}$. 

 Doing some algebra on the top and bottom we get:

 \begin{equation}2^a \cdot \frac{1+2^{k(m+1)}}{1+2^k} = 964\end{equation}

 The prime factorization of $964$ is $241 \cdot 2^2$.

 Equating $2^a = 2^2$, we get $a = 2$.

 Next since $241$ is a prime number we equate the second term of the product to $241$.

 \begin{equation}\frac{1+2^{k(m+1)}}{1+2^k} = 241\end{equation}Doing some algebra we get, that $2^k(2^{km}-241) = 240$

 The closest power of $2$ to $241$ is $256$ which is $2^8$. So setting $km = 8$, we get $2^8-241 = 15$
$2^k(15) = 240$, $2^k = 16$, $k = 4$.

 This we know that if $k = 4$, then $4m = 8$, so $m = 2$.

 Finally, we conclude that $a = 2$, $m = 2$, and $k = 4$.

 
$4+2+2 = 8$, so the answer is $\fbox{A}$

 Actually when it comes to $2^k(2^{km}-241) = 240$;we can directly observe the answer by discovering $2^k$ is even and $2^{km}-241$ must be odd. So, the only available divisors of 240 is 16 and 15.And we can have $m = 2$, and $k = 4$.

 



\textbf{Solution 2(alternative method)}

We start from this equation derived in Solution 1:

 \begin{equation}\frac{1+2^{k(m+1)}}{1+2^k} = 241\end{equation}

 Note that $241=256-16+1=16^2-16+1$. Also note that this looks similar to the second term in the Sum of Cubes factorization $a^3+b^3=(a+b)(a^2-ab+b^2)$. Specifically, when $a=16$ and $b=1$, the right term is $16^2-16+1$. Thus, $\frac{16^3+1^3}{16+1}=241$. We rewrite: $\frac{1+2^{12}}{1+2^4}=241$

 We can then get that $m = 2$ and $k=4$, and proceed as above.

 


	\item (\textit{AMC 12 12A 2025}) Rectangular card $ABCD$ has $AB = 1$ and $AC = 2$. As shown in the figure, a second card, $AEFG$, congruent to the first, is stacked onto it so that $\overline{AG}$ lies on $\overline{AC}$ and the two cards share corner $A$.


 
                    \begin{figure}[H]
                        \centering
                        \includegraphics[width=0.5\textwidth]{Problem/AMC_12_12A_2025_22_Problem_Figure_1.png}
                    \end{figure}
                    


 Additional cards are added, following the same pattern (with shared corner $A$), creating a fan of cards. The process stops when a new card fits exactly atop $ABCD$. The shape then created is a nonconvex polygon with rotational symmetry. What is the area of this polygon?


 $\textbf{(A)}\   20\sqrt{3} - 24  \qquad                 \textbf{(B)}\   34\sqrt{3} - 48  \qquad                 \textbf{(C)}\   6 + 3\sqrt{3}  \qquad                 \textbf{(D)}\   12  \qquad                 \textbf{(E)}\   8\sqrt{3}$

 \textbf{Answer: }E



\textbf{Solution 1 (3D geometry)}

Let the three numbers be $X$, $Y$, and $Z$, and consider each possible triple $(X,Y,Z)$ as the (Cartesian) coordinates of a point in 3-dimensional space. Then, as each number is chosen independently and uniformly from the interval $[0,1]$, the total space of possible points $(X,Y,Z)$ is the unit cube $0 \leq X,Y,Z \leq 1$. Using the idea of \href{https://artofproblemsolving.com/wiki/index.php?title=Geometric\_probability}{geometric probability}, it follows that the required probability equals the volume of the region consisting of all points $(X,Y,Z)$ that satisfy the required condition, divided by the volume of the total space, i.e. the unit cube (which is simply $1^3 = 1$). Accordingly, it only remains to determine the region of possible points $(X,Y,Z)$ and then compute its volume.

 For simplicity, we may assume that e.g. $X$ is the maximum - clearly, by symmetry, the corresponding regions where $Y$ is the maximum and where $Z$ is the maximum have the same volume, and these regions overlap only where at least two of $X,Y,Z$ are both the maximum. In such cases, those two variables must be equal (e.g. $X$ and $Y$ will both be the maximum precisely if $X = Y \geq Z$), and so each of these overlaps is either a region within a plane (as e.g. the equation $X = Y$ defines a plane), or a straight line (in the special case where $X = Y = Z$). It follows that the overlaps between the $3$ regions all have zero volume, so since each region has the same volume (as explained above), the answer will simply be $3$ times the volume of one of these regions, say the one where $X$ is the maximum and the required condition is satisfied.

 Within that region, we therefore have $0 \leq X,Y,Z \leq 1$, $X \geq Y$ and $X \geq Z$ for $X$ to be the maximum, and $X > 2Y$ and $X > 2Z$ for the required condition, but this combination of inequalities reduces to simply $0 \leq 2Y,2Z < X \leq 1$, since we automatically have $2Y \geq Y$ and $2Z \geq Z$ as $Y,Z \geq 0$. This means the boundaries of the region are the lines $X = 2Y$, $X = 2Z$, $X = 0$, $Y = 0$, and $X = 1$, so now by drawing the unit cube as well as those boundaries, we deduce that this region is a square-based pyramid: the base has vertices $\left(1,0,0\right)$, $\left(1,0,\frac{1}{2}\right)$, $\left(1,\frac{1}{2},\frac{1}{2}\right)$, and $\left(1,\frac{1}{2},0\right)$, while the apex of the pyramid is at $(0,0,0)$. The square base thus clearly has side length $\frac{1}{2}$, and hence area $\left(\frac{1}{2}\right)^2 = \frac{1}{4}$; moreover, as the base lies in the plane $X = 1$, the perpendicular height of this pyramid is the same as the perpendicular distance from the apex $(0,0,0)$ to that plane, which is obviously $1-0 = 1$.

 It follows that the volume of this pyramid, i.e. of the region where $X$ is the maximum and the required condition is satisfied, is $\frac{1}{3} \cdot \frac{1}{4} \cdot 1 = \frac{1}{12}$ (using the standard formula for the volume of a pyramid). Finally, as explained above, this means the required probability is $3 \cdot \frac{1}{12} = \boxed{\textbf{(E)}~\frac14}$.

 



\textbf{Solution 2 (calculus)}

Consider three independent uniform random variables $X, Y, Z$ on $[0, 1]$. The goal is to find the probability that the maximum value is greater than twice each of the other two values. Equivalently, if we order them as $a \geq b \geq c$, this is the probability that $a > 2b$ (which implies $a > 2c$ since $b \geq c$). Therefore, if we assume that $X$ is the maximum (but make no assumption about the order of $Y$ and $Z$), the required condition reduces (almost surely) to $X > 2Y$ and $X > 2Z$. Conversely, if $X > 2Y$ and $X > 2Z$, then we will necessarily have $X > 2Y \geq Y$ and $X > 2Z \geq Z$ since $Y,Z \geq 0$, and so $X$ will indeed be the maximum.

 Accordingly, the probability that $X$ is the maximum and the required condition is satisfied is simply $P(X > 2Y \cap X > 2Z)$, and then the answer will be $3$ times this probability, to account for the $3$ total cases (the maximum can be $X$, $Y$, or $Z$, and clearly all of these cases are symmetric). Since $X,Y,Z$ are chosen uniformly at random from the interval $[0,1]$, the marginal distribution of each one is simply the uniform distribution on this interval. As they are also chosen independently, we deduce that their joint distribution is uniform over $[0,1] \times [0,1] \times [0,1]$, i.e. the unit cube, and hence the joint probability density is $\frac{1}{\text{volume of the unit cube}} = \frac{1}{1^3} = 1$.

 
For the bounds of integration, we can first choose $X \in [0,1]$, and then we require $Y,Z \in [0,1]$, $X > 2Y$, and $X > 2Z$; as $X \leq 1$, the latter two conditions will also automatically give $Y,Z < \frac{X}{2} \leq \frac{1}{2} < 1$, so it is both necessary and sufficient that $0 \leq Y,Z < \frac{X}{2}$. Thus we can now compute
 \begin{align*} P(X > 2Y \cap X > 2Z) &= \int_{0}^{1} \int_{0}^{\frac{x}{2}} \int_{0}^{\frac{x}{2}} 1 \, \mathop{}\!\mathrm{d}z \, \mathop{}\!\mathrm{d}y \, \mathop{}\!\mathrm{d}x \\ &= \int_{0}^{1} \left(\frac{x}{2}\right)^2 \, \mathop{}\!\mathrm{d}x \\ &= \int_{0}^{1} \frac{x^2}{4} \, \mathop{}\!\mathrm{d}x \\ &= \frac{1}{4}\left[\frac{1}{3}x^3\right]_{0}^{1} \\ &= \frac{1}{4} \cdot \frac{1}{3} \left(1^3-0^3\right) \\ &= \frac{1}{12}, \end{align*}
 and so, as explained above, the total probability is $3 \cdot \frac{1}{12} = \boxed{\textbf{(E)}~\frac{1}{4}}$.

 



\textbf{Solution 2.1 (alternative, fixing the entire order)}

Alternatively, instead of just fixing the maximum as $X$, we can fix the entire order of the $3$ variables as e.g. $X > Y > Z$. It then follows that the required condition reduces (almost surely) to $X > 2Y > 2Z$, while conversely, if $X > 2Y > 2Z$, it will indeed necessarily follow that the order is $X > Y > Z$, as $Y \geq 0$ again gives $X > 2Y \geq Y$ and $2Y > 2Z$ gives $Y > Z$.

 Hence, similarly to above, the probability that the required condition is satisfied is $P(X > 2Y > 2Z)$, and then as $X > Y > Z$ is just $1$ of the total $3! = 3 \cdot 2 \cdot 1 = 6$ possible orders (with all of them clearly being symmetric), the answer will be $6$ times this probability. The joint probability density is still $1$ as shown above, but now for the bounds of integration, we first choose $Y$, where we must have $1 \geq X > 2Y$, and thus $0 < Y < \frac{1}{2}$. We can now choose $X$, again satisfying $X > 2Y$ and of course $0 \leq X \leq 1$, i.e. $2Y < X \leq 1$ (since we then automatically have $X > 2Y \geq 2 \cdot 0 = 0$). Lastly, we need to choose $Z$ such that $Y > Z$ and $0 \leq Z \leq 1$ (as then $X > 2Y$ will automatically also give $X > 2Z$), i.e. $0 \leq Z < Y$ (since $Y \leq 1$ will then automatically give $Z \leq 1$).

 Accordingly, we obtain
 \begin{align*} P(X > 2Y > 2Z) &= \int_{0}^{\frac{1}{2}} \int_{2y}^{1} \int_{0}^{y} 1 \, \mathop{}\!\mathrm{d}z \, \mathop{}\!\mathrm{d}x \, \mathop{}\!\mathrm{d}y \\ &= \int_{0}^{\frac{1}{2}}y(1-2y) \, \mathop{}\!\mathrm{d}y \\ &= \int_{0}^{\frac{1}{2}}\left(y-2y^2\right) \, \mathop{}\!\mathrm{d}y \\ &= \left[\frac{1}{2}y^2-\frac{2}{3}y^3\right]_{0}^{\frac{1}{2}} \\ &= \left(\frac{1}{2} \cdot \left(\frac{1}{2}\right)^2 - \frac{2}{3} \cdot \left(\frac{1}{2}\right)^3\right) - (0-0) \\ &= \frac{1}{24}, \end{align*}
and so, as explained above (and just as obtained in Solution 2), the required probability is $6 \cdot \frac{1}{24} = \boxed{\textbf{(E)}~\frac{1}{4}}$.

 



\textbf{Solution 3 (easiest and fastest and simplest and bestest)}

WLOG, let $a$ be the greatest of the three numbers. The probability that it is twice as great as one of the other numbers chosen is $\frac{1}{2}$. This is because there is an equal probability that the number is chosen from $(0,\frac{a}{2})$ or $(\frac{a}{2},a)$. So the probability that both $b$ and $c$ are less than half of $a$ is $\frac{1}{2} \cdot \frac{1}{2} = \boxed{\textbf{(E)}~\frac{1}{4}}$.

 



\textbf{Solution 3}

Let $a,b,c$ be three independent random variables distributed uniformly on $[0,1]$. By symmetry, either 3 could be the greatest. WLOG, let $a$ be the greatest.When $a=x$, we need to have $b,c<x/2$. Since $b$ and $c$ are uniformly distributed, the probability is $(x/2)^2=x^2/4$. Integrate:
 \[\int_0^1\frac{x^2}{4}dx=\frac{1}{12}\]Multiply it by 3 to account for symmetry:
 \[\frac{1}{12}\cdot3=\boxed{\frac{1}{4}}\]~dondee123

 


	\item (\textit{AMC 12 12A 2025}) For positive integers $a$ and $b$, let $R(a, b)$ denote the remainder when $a$ is divided by $b$. Suppose $N$ is the least positive integer such that

 \[R(N,2) \ge 1, \quad R(N,4) \ge 2, \quad R(N,8) \ge 4, \quad \dots, \quad R\left(N,2^{200}\right) \ge 2^{199}.\]
What are the last three digits of $N$?


 $\textbf{(A)}\   125  \qquad                 \textbf{(B)}\   249  \qquad                 \textbf{(C)}\   375  \qquad                 \textbf{(D)}\   511  \qquad                 \textbf{(E)}\   875$

 \textbf{Answer: }C



\textbf{Solution 1}

To satisfy the conditions, a fair integer must have no digit be a local minimum. Let's say we have $n$ distinct digits, with each digit being a number from $1$ to $9$. To create a fair integer, we begin by placing the largest digit. For the second-largest digit, we can either place this digit to the right or to the left of the string already created. We have these $2$ options for the third-largest digit, and so on. Therefore, there are $2^{n-1}$ valid permutations to create a fair integer.

 We must also choose which digits will be in the permutation. If you are creating an $n$-digit long fair integer, there are $\binom{9}{n}$ ways to pick which digits will be in the number.

 Therefore, for each $n \in \{1,2,\dots, 9\}$, the number of fair integers of length $n$ is:
 \[\binom{9}{n} \cdot 2^{n-1}.\]Summing over all $n$:
 \[\sum_{n=1}^9{\binom{9}{n} \cdot 2^{n-1}}=\frac{1}{2}\left(\sum_{n=0}^9{\binom{9}{n}}2^n -1\right)=\frac{1}{2}\left((1+2)^9 -1 \right) = \frac{1}{2}(19682) = \boxed{9841}.\]Note that the Binomial Theorem was used to equate
 \[\sum_{n=0}^9{\binom{9}{n}}2^n = (1+2)^9.\]~lprado

 



\textbf{Solution 2}

Note every fair number will have an increasing string of digits, a maximum digit, then a decreasing string of digits. This is because if it decreases then increases, then the digit in the middle will be less than its adjacent digits.

 Let $n$ be the maximum digit. For each number $i<n$, we may either place $i$ before $n$, after $n$, or choose not to include it. Note this process will result in a unique number for every case, as the numbers before $n$ must be in increasing order, and the numbers after $n$ must be in decreasing order. Therefore, for each number $n$, we have $3^{n-1}$ cases.

 Since $n \in \{1,2,\cdots9\}$, we have:
 \[\sum_{n=1}^9 3^{n-1}=\frac{3^9-1}{3-1}=\boxed{9841}.\]




\textbf{Solution 2.5 (Cheese, Goldilocks Zone)}

If you consider options $D$ and $E$ too big, and $A$ and $B$ too small, then $\boxed {\text {(C) } 9841}$ is your answer.

 Also $E$ is twice $C$ so $C$ is most likely your answer (since for double counting etc.). 

 - Ameen Patel IrvingtonHighSchoolMathGeeks-Edited by AlgeBruh16-Edited by TFEA

 



\textbf{Solution 3 (Recursion)}

Let $f(n)$ be the number of fair integers given an arbitrary set of $n$ digits. Let $a$ be the smallest of these digits. Notice that $a$ is either the first or last digit, as if it were any other digit, the two digits surrounding it would both be greater. Then, notice that if the remaining $n-1$ slots form a fair number, so does the number when $a$ is appended. Therefore, $f(n) = 2f(n-1)$. From here, we may proceed with the calculation in Solution $1$ to get the answer of $\boxed {\text {(C) } 9841}$




\textbf{Solution 4}

For any fair integer $n$, we write some $9$-digit base-$3$ number(s) $f(n)=\left(\overline{d_1 d_2 \ldots d_9}\right)_3$ from left to right as follows:

 
            \begin{itemize}
        \item $d_k=0$ if $k$ is not in $n$;
        \item $d_k=1$ if $k$ is the to the left of all digits greater than $k$ in $n$;
        \item $d_k=2$ if $k$ is the to the right of all digits greater than $k$ in $n$.
            \end{itemize}
            Note that the fairness ensured one and only one of the above three cases happens. 

 Now if $k$ is the largest digit in $n$ then $d_k$ can be either $1$ or $2$.

 An example: If $n=136852$, then $f(n)=\overline{121021010}_3$ or $\overline{121021020}_3$. 

 It's thus easy to see that there is a $1$-to-$2$ mapping from all fair integers to all $9$-digit base-$3$ nonzero numbers. Therefore the answer is $\frac{3^9-1}{2}$.

 



\textbf{Solution 5 (similar to [3])}

\textbf{Lemma:} Consider any $n$ digit fair integer, and drop the smallest digit. The $n-1$ version of a fair integer should be a fair integer.

 \textbf{Proof of the Lemma:}

 If $\overline{a_1 a_2 a_3 a_4 \dots a_k}$ with $k \leq 9$ is the integer, then firstly notice that if $a_1$ is the smallest, then $a_1$ must be on either one of the endpoints. If you drop this digit, $\overline{a_2 a_3 a_4 \dots a_k}$ is an integer. Since adding $a_1$ to either one of the endpoints is an identity (as $a_1$ is the smallest), the given integer must be fair as otherwise its $n$ digit interpretation is not fair. $\blacksquare$.

 \textbf{2025 AMC 12A P23:}

 From an $n-1$ string, we can choose a smaller digit and place it in two ways, the start or end point, and we have the recurrence, if $f(n)$ denotes the number of fair integers given $n$ digits: $f(n) = 2f(n-1)$. 

 One can simply bash the recurrence from here to obtain $\boxed{9841}$.

 


	\item (\textit{AMC 12 12A 2025}) Let $\triangle ABC$ be a scalene triangle with incenter $I$. (Recall that the incenter is the intersection of the three angle bisectors.) The distances from $I$ to the midpoints of the sides of $\triangle ABC$ are $\sqrt{11}$, $\sqrt{14}$, and $\sqrt{19}$. What is the sum of the possible areas of the circle inscribed in $\triangle ABC$?


 $\textbf{(A)}\   10 \pi  \qquad                 \textbf{(B)}\   \frac{44 \pi}{3}  \qquad                 \textbf{(C)}\   \frac{58 \pi}{3}  \qquad                 \textbf{(D)}\   \frac{88 \pi}{3}  \qquad                 \textbf{(E)}\   30 \pi$

 \textbf{Answer: }C



\textbf{Solution 1 (Sin 15)}

Let the center of the large circle be $O$ and the centers of the $12$ circles be $A_1, A_2, A_3, \dots, A_{12}$. Triangle $OA_1A_2$ has side lengths $r+1, r+1, 2$, with the angle opposite $2$ being $360/12 = 30$. Note that the line connecting $A_1$ and $A_2$ go through their common point of tangency, by definition, which causes $A_1A_2$ to have a length of $2$.

 Drawing the angle bisector of the $30$ degree angle, we split $OA_1A_2$ into two congruent right triangles, each with hypotenuse $r+1$ and side opposite the $15$ degree angle $1$.

 From here, note that $\sin{15} = \frac{\sqrt{6}-\sqrt{2}}{4}$, which be derived using the trigonometric identity $\sin{(A-B)} = \sin{A} \cos{B} - \sin{B} \cos{A}$, with $A=45$ and $B=30$.

 In our right triangle, $\sin{15} = \frac{1}{r+1} = \frac{\sqrt{6}-\sqrt{2}}{4}$. Let $x=r+1$. Solving for $x$, we get $x = \frac{4}{\sqrt{6}-\sqrt{2}}$. Rationalizing, we get that $x = \sqrt{6}+\sqrt{2}$.

 Remember $x = r+1 = \sqrt{6}+\sqrt{2}$, so $r = \sqrt{6}+\sqrt{2} - 1$. Therefore, our answer is $6+2-1 = \boxed{7}.$




\textbf{Solution 2 (Law of Cosines Bash)}

Let the center of the large circle be $O$ and the centers of any two circles be $A$ and $B$. Triangle $OAB$ has side lengths $r+1, r+1, 2$, with the angle opposite $2$ being $360/12 = 30$.

 Using Law of Cosines, $2^2=AO^2+BO^2-2AOBO\cos30$. Plugging in the radius yields $2^2=(r+1)^2+(r+1)^2-2(r+1)^2\cdot\frac{\sqrt{3}}{2}$. After solving for $r$ and simplifying, we get $r=\sqrt6+\sqrt2-1$. Therefore our answer is $7$.

 \textbf{Worked Out Solution} $4=(2-\sqrt{3})(r+1)^2$. Divide both sides to get $\frac{4}{2-\sqrt{3}}=(r+1)^2$. Rationalize and expand to get $8+4\sqrt{3}=r^2+2r+1$. Then $r^2+2r+(-7-4\sqrt{3})=0$. Use the quadratic formula to get $r=\frac{-2+\sqrt{4+28+16\sqrt{3}}}{2}$, or $r=-1+\sqrt{8+4\sqrt{3}}$ (Note: We only take the plus instead of plus-minus because r must be positive). Clearly $c=-1$, and then we need to find some a and b such that $\sqrt{a} + \sqrt{b} = \sqrt{8+4\sqrt{3}}$. If we square both sides, $a+b+2\sqrt{ab} = 8+4\sqrt{3}$, so $a+b=8$. Plug in our previously found values to find $a+b+c=8+(-1)$. Therefore our answer is $7$.

 
$\textbf{Remark}$

 We could've also noticed that $\sqrt{2-\sqrt{3}} = \sqrt{\frac{1}{2}\cdot(4-2\sqrt{3})} = \frac{1}{\sqrt{2}}*\sqrt{3-2\sqrt{3}+1} = \frac{1}{\sqrt{2}}*(\sqrt{3}-1)$, and therefore directly gotten $r+1 = \sqrt{\frac{4}{2-\sqrt{3}}} = \frac{2}{\frac{1}{\sqrt{2}}*(\sqrt{3}-1)} = \frac{2\sqrt{2}}{\sqrt{3}-1}$ and solved from there.

 $\textbf{Commentary}$

 This solution is so good but it can be frustrating when such an important detail is glossed over. "Clearly $c=-1.$" While this seems harmless, it is simple to not know to make this assumption and get stuck with too many variables. I would be curious if someone could come up with a rigorous response to this.

 



\textbf{Solution 3 (Easy Geo Method)}

We connect the centers of the smaller circles at $0^\circ$ and $90^\circ$ as follows:
                    \begin{figure}[H]
                        \centering
                        \includegraphics[width=0.5\textwidth]{Solution/AMC_12_12A_2025_24_Solution_Figure_1.png}
                    \end{figure}
                    

 Note that connecting the centers of the 12 smaller circles forms a regular 12-gon. A regular 12-gon has angles 

 
 \[\frac{10\cdot 180}{12} = 150^\circ,\] 

 so $\angle OAB = \angle ODC = 75^\circ.$ Furthermore, since $OA=OD$ and $\angle AOD = 90^\circ$, we have $\angle OAD = 45^\circ$, so $\angle DAB = 30^\circ$. Next, we drop the perpendiculars from $B$ and $C$ to $AD$. This creates a $1 \times 2$ rectangle and two $30-60-90$ triangles. Thus,

 
 \[AD = (r+1)\sqrt{2} = 2+2\sqrt{3},\]

 and manipulating yields

 
 \[r = -1 + \sqrt{2}+\sqrt{6},\]

 so the desired answer is $-1+2+6 = \boxed{\text{(C) } 7.}$

 P.S. I posted the asymptote code here. Anyone that wishes to use it can steal~ :)

 



\textbf{Solution 4 (Fast Law of Sines)}

Add a super simple construction to the diagram as shown below (i.e. partition 1/12th of the dodecagon whose vertices are the centers of the circles)
                    \begin{figure}[H]
                        \centering
                        \includegraphics[width=0.5\textwidth]{Solution/AMC_12_12A_2025_24_Solution_Figure_2.png}
                    \end{figure}
                    Asymptote's being difficult with me (or rather I just suck at it) so see that we have $\angle OBA=\angle OAB=75^\circ$ and $\angle BOA=30^\circ.$ Using the law of sines, we have
 \[\frac{r+1}{\sin(75^\circ)}=\frac{2}{\sin(30^\circ)}=4\implies r=4\sin(75^\circ)-1.\]We now just compute
 \[\sin(75^\circ)=\sin(30^\circ+45^\circ)=\sin(30^\circ)\cos(45^\circ)+\cos(30^\circ)\sin(45^\circ)=\frac{\sqrt 3+1}{2\sqrt{2}}=\frac{\sqrt 6+\sqrt 2}{4}.\]Thus
 \[r=\sqrt 6+\sqrt 2-1\implies a+b+c=\boxed{7}.\]~shockfront99 (asy stolen from ABC09090927)

 


	\item (\textit{AMC 12 12A 2025}) A \textit{finite integer tower} is an expression of the form

 \[a_{1}^{a_{2}^{a_{3}^{\iddots^{a_{n}}}}}\]
where $n \geq 1$ and $a_{1}$, $a_2$, $\dots$, $a_{n}$ are positive integers greater than $1$. (Recall that exponentiation is performed from the top down. For example, $3^{2^{3}}$ equals $3^{8}$, not $9^{3}$.) Two towers are considered different if they have a different number of layers (values of $n$) or at least one pair of numbers in corresponding positions differ.


 Define $f(x)$ to be the number of different finite integer towers whose value is $x$. (For example, $f(16)=4$ because $4$ towers evaluate to $16$, namely $16$, $2^{4}$, $4^2$, and $2^{2^{2}}$.) What is $f(2^{8192})$?


 $\textbf{(A)}\   17  \qquad                 \textbf{(B)}\   32  \qquad                 \textbf{(C)}\   42  \qquad                 \textbf{(D)}\   43  \qquad                 \textbf{(E)}\   47$

 \textbf{Answer: }None of the official answer choices were correct.



\textbf{Solution 1 (Got Wrong)}

None of the answer choices on the official test (which asked for the number of possible functions $f(x)$) were correct, but choice E would be correct if this problem asked for the number of pairs of functions $(P(x), Q(x))$.

 Let $R(x) = \frac{P(x)}{Q(x)}$.  Since $R(x) \leq 0$ on $[a, b]$ but not for values slightly less than $a$ or slightly more than $b$, $P(x) = 0$ at $x = a$ and $x = b$.  Therefore, $P(x) = (x-a)(x-b)(x-r)$ for some $r \in \{1, 2, 3, 4, 5\}$.

 Since $R(x) \leq 0$ on $(c, d)$ but not at $x = c$ or $x = d$, $R(x)$ is not continuous at $x = c$ or $x = d$.  Therefore, $R(x)$ must be undefined at $x = c$ and $x = d$, that is, $Q(x) = 0$ at $x = c$ and $x = d$.  So $Q(x) = (x-c)(x-d)(x-s)$ for some $s \in \{1, 2, 3, 4, 5\}$. 

 Therefore, $R(x) = \frac{(x-a)(x-b)(x-r)}{(x-c)(x-d)(x-s)}$.  Notice that $\frac{(x-a)(x-b)}{(x-c)(x-d)} \leq 0$ only on $[a, b]$ and $(c, d)$.  Therefore, $\frac{x-r}{x-s}$ must be nonnegative for all $x \notin \{a, b, c, d, r, s\}$.  This only happens if $r = s$.

 Thus $R(x) = \frac{(x-a)(x-b)(x-r)}{(x-c)(x-d)(x-r)}$, which is the same as $\frac{(x-a)(x-b)}{(x-c)(x-d)}$ except that it is undefined at $x = r$.  Thus $R(x)$ satisfies the desired property as long as $r \notin [a, b] \cup (c, d)$.

 Note that each quintuple $(a, b, c, d, r)$ defines a unique pair of functions $(P(x), Q(x))$.

 If $(a, b, c, d) = (1, 2, 3, 4)$, $r$ can be $3$, $4$, or $5$.

 If $(a, b, c, d) = (1, 2, 3, 5)$, $r$ can be $3$ or $5$.

 If $(a, b, c, d) = (1, 2, 4, 5)$, $r$ can be $3$, $4$, or $5$.

 If $(a, b, c, d) = (1, 3, 4, 5)$, $r$ can be $4$ or $5$.

 If $(a, b, c, d) = (2, 3, 4, 5)$, $r$ can be $1$, $4$, or $5$.

 Therefore, there are $\boxed{(\textbf{E})\ 13}$ possible pairs of functions $(P(x), Q(x))$.

 -j314andrews

 



\textbf{Solution 2 (credit to Sohil Rathi's video solution, pi_is_3.14 / OmegaLearn)}

This is the solution to the original problem, which asked for the number of possible \textbf{functions}, not pairs of polynomials. 

 First, consider the problem if it were talking about two second degree polynomials. We can see that the function 
 \[f(x) = \frac{(x-a)(x-b)}{(x-c)(x-d)}\] by itself satisfies the condition of dropping below the x axis over $[a,b] \cup (c,d).$ Now, we need to add one extra $(x-n)$ term each to the numerator and denominator. 

 Case 1: $f(x) = \frac{(x-a)(x-b)(x-f)}{(x-c)(x-d)(x-f)}$ for some other value $f$ not equal to $a, b, c,$ or $d.$ 

 Note that $f$ cannot be between $a$ and $b$ or between $c$ and $d$ because this would create a hole in the interval. So the only possibilities are: 

 
            \begin{itemize}
        \item $a = 1, b = 2, c = 3, d = 4, f = 5$
        \item $a = 1, b = 2, c = 4, d = 5, f = 3$
        \item $a = 2, b = 3, c = 4, d = 5, f = 1$
            \end{itemize}
            This gives us $3$ valid functions so far.

 
Case 2: $f(x) = \frac{(x-a)(x-b)(x-f)}{(x-c)(x-d)(x-f)}$ where $f$ is $c$ or $d.$

 If $f$ equals $c$ or $d$ the resulting function is essentially equivalent to $f(x) = \frac{(x-a)(x-b)}{(x-c)(x-d)}$ because a hole would already exist, so we can safely cancel the $(x-f)$ terms. It doesn't matter whether we choose $f$ to equal $c$ or $d$ because the function will still be the same. There are $\binom{5}{4} = 5$ ways to select the values for $a, b, c,$ and $d$ and $1$ way to choose any variable for $f$ to equal, giving $5 \cdot 1 = 5$ valid functions for this case. 

 
The correct answer is $5 + 3 = \boxed{\textbf{(F) } 8}$ functions.

 
i tried my best to show what I think was Sohil's thought process, but please feel free to edit this if you can explain it better :)

 



\textbf{Solution 3}

\textbf{Solution 1} by user j314andrews is incomplete; \textbf{Solution 2} by user \href{https://artofproblemsolving.com/wiki/index.php?title=User:Grogg007}{grogg007} commits a serious logical error because it only provides functions that satisfy the requirements, without showing that no other functions do. Other places like \href{https://wiki.randommath.com/amc12/2025/part-a/problem-25}{[1]} gives a good solution, but unfortunately it requires calculus. Here I am trying to provide a good solution that does not require calculus.

 
Let $D$ and $E$ be the set of the real roots of $P(x)$ and $Q(x)$, respectively. Then, outside of the intervals $[a,b]$ and $(c,d)$, $f(x)>0$ except when $x \in E$, where $f(x)$ is undefined. 

 
For any function $g(x)$, define a real number $t$ as a \textbf{toggling point} for $g(x)$ if there exists a sequence $t_1 < t_2 < ... < t_N$, with the understanding that $t_0 = -\infty$ and $t_{N+1} = \infty$, such that 

 
            \begin{itemize}
        \item $t = t_k$ for some $1 \leq k \leq N$,
        \item $g(x)$ is of the same sign (either always positive or always negative) on each interval $(t_j, t_{j+1})$, for $j=0, ..., N$, and
        \item $g(x)$ has opposite signs on $(t_{k-1}, t_k)$ and $(t_k, t_{k+1})$.
            \end{itemize}
            We call these open intervals $(t_j, t_{j+1})$ ($j=0, ..., N$) the \textbf{bands} of $g(x)$. 

 
The polynomial $P(x)$ has no toggling point outside of $D$, and $Q(x)$ has no toggling point outside of $E$. Now in our problem, $f(x)$ has exactly 4 toggling points: $a, b, c, d$. Our task is to see how these four points are related to the sets $D$ and $E$.  

 \textbf{Lemma}. At each toggling point $t$ for $f(x)$, 

 
            \begin{enumerate}
        \item either $P(t)=0$ or $Q(t)=0$ (or both);
        \item if the multiplicities of the root $t$ of $P(x)$ and $Q(x)$ are $m$ and $n$, respectively, then $m-n$ is odd.
            \end{enumerate}
            \textbf{Proof}. (1) if both $P(t)$ and $Q(t)$ are nonzero, then $t$ belongs to a band of $P(x)$ and a band of $Q(x)$. Taking the intersection between these two bands, we get an open interval on which $f(x)$ has the same sign as $f(t)$. Thus $t$ cannot be a toggling point of $f(x)$. 

 (2) If $m$ is odd, there are real numbers $r, s$, with $r<t<s$, such that $P(x)$ has the same sign on $(r,t)$ and has the opposite sign on $(t, s)$. If $m$ is even, there are real numbers $r, s$, with $r<t<s$, such that $P(x)$ has the same sign on $(r,t)$ and on $(t,s)$. Similar statements can be made of $Q(x)$, having intervals $(u, t)$ and $(t, v)$ so that $Q(t)$ is of the same sign on $(u,t)$, and of either a different or the same sign on $(t, v)$, depending on the parity of $n$.  If both $m$ and $n$ are odd or both $m$ and $n$ are even, then $P(t)/Q(t)$ is of the same sign on the interval $(\max\{r, u\}, t)$ and $(t, \min\{s, v\})$, so $t$ cannot be a toggling point for $f(x)$.

 
\textbf{Solution to the problem}. Since $a$ is a toggling point for $f(x)$, either $P(a)=0$ or $Q(a)=0$, or both. But since $f(a)$ is defined, we have $Q(a)\neq 0$, so $P(a)=0$. The same can be said of $b$, so $P(b)=0$. The multiplicities of $a$ and $b$ must be $1$. Otherwise, point (2) of the Lemma would require the multiplicity of one of them to be at least $3$, but $P(x)$ is only cubic.

 
$c$ is a toggling point for $f(x)$, so either $P(c)=0$ or $Q(c)=0$, or both. But if $Q(c)\neq 0$, we have $f(c)=0$, violating the condition. So $Q(c)=0$. Similarly, $Q(d)=0$. Concerning their multiplicities, point (2) of the Lemma gives us three cases, given that $Q(x)$ is only cubic: 

 
            \begin{itemize}
        \item (a) $c$ is a single root of $Q(x)$, $P(c)\neq 0$, $d$ is a double root of $Q(x)$ and a single root of $P(x)$;
        \item (b) switching the roles of $c$ and $d$;
        \item (c) both $c$ and $d$ are single roots of $Q(x)$, and $Q(x)$ has a single root $e$ other than $c$ and $d$.
            \end{itemize}
            
Case (a) provides functions in the form of $\frac{(x-a)(x-b)(x-d)}{(x-c)(x-d)^2}$. There are 5 of them because there are 5 ways to choose $a, b, c, d$ from the set $\{1, 2, 3, 4, 5\}$.

 Case (b) provides functions in the form of $\frac{(x-a)(x-b)(x-c)}{(x-c)^2(x-d)}$, but they are the same as those provided in case (a). To ensure this, notice that the two expressions have the same domain and they are equal for every $x$ in their domains.

 Case (c) requires that $P(e)=0$, because if $P(e)\neq 0$ then $e$ becomes an extra toggling point for $f(x)$. So this case provides functions in the form of $\frac{(x-a)(x-b)(x-e)}{(x-c)(x-d)(x-e)}$. But since $f(e)$ is undefined, we need to make sure that $e$ does not fall into $[a,b]$ or $(c,d)$. So $e$ can be $1, 3, 5$. In all these cases, we also have $f(e) = \frac{(e-a)(e-b)}{(e-c)(e-d)} > 0$, satisfying the conditions. That's 3 more functions.

 So the total number of functions is $5 + 3 = \boxed{8}$.

 \href{https://artofproblemsolving.com/wiki/index.php?title=User:Lightest}{Lightest} (\href{https://artofproblemsolving.com/wiki/index.php?title=User\_talk:Lightest&action=edit&redlink=1}{talk}) 22:51, 10 November 2025 (EST)

 



\textbf{Solution 4 (Multiplicity Analysis)}

For the condition $f(x)\le 0 \iff x\in [a,b]\cup(c,d)$, $f$ must change sign at $a, b, c,$ and $d$. The key trick is that a rational function can only change sign through zeroes or discontinuities. Since both $P(x)$ and $Q(x)$ are monic cubic polynomials, $f(x) \to 1$ as $x \to \pm\infty$. This means $f(x)$ is $+$ for large $|x|$. The required sign pattern must therefore be:
 \[f(x): \quad +\underset{ a }{ [ }-\underset{ b }{ ] }+\underset{ c }{ ( }-\underset{ d }{ ) }+\]

            \begin{itemize}
        \item For the interval to be closed, $[a, b]$, we must have $f(a) = 0$ and $f(b) = 0$. This implies $P(a) = P(b) = 0$ and, critically, $Q(a) \neq 0$ and $Q(b) \neq 0$ to avoid discontinuity.
        \item For the interval to be open, $(c, d)$, the function must be undefined at the endpoints: if $f(c)$ were defined, by the intermediate value theorem $f(c)=0$, violating the condition. This implies $Q(c) = 0$ and $Q(d) = 0$. Note that $P(c),P(d)$ could be $0$, as $\frac{0}{0}$ remains undefined.
            \end{itemize}
            From this, we conclude that $a, b, c, d\in \{ 1,2,3,4,5 \}$. Let the remaining root in be $t$, so $\{ 1,2,3,4,5 \}=\{ a,b,c,d,t \}$.

 Let $p_i, q_i\in \mathbb{Z}_{\ge 0}$ be the multiplicities of a root $i=1,2,3,4,5$ in $P(x),Q(x)$, respectively, e.g. $P(x)=(x-1)^{p_{1}}(x-2)^{p_{2}}(x-3)^{p_{3}}(x-4)^{p_{4}}(x-5)^{p_{5}}$. Let the multiplicity in $f(x)$ at $i$ be $m_i = p_i - q_i$. A sign change occurs at $i$ if and only if $m_i$ is odd.

 
            \begin{itemize}
        \item The sign changes at $a, b, c, d$. Thus, $m_a, m_b, m_c, m_d$ must all be odd.
        \item The sign does not change at $t$. Thus, $m_t$ must be even.
            \end{itemize}
            From the root analysis, $q_a = 0 \implies m_a=p_{a}-q_{a} = p_a$ is odd. Similarly, $p_b$ is odd. Since $\deg(P) = \sum p_i = 3$, and $p_a, p_b\ge 1$ are odd, $m_{a}=m_{b}=p_{a}=p_{b}=1$. This forces the roots of $P(x)$ to be distinct: $\{a, b, k\}$ for some $k \in \{ 1,2,3,4,5 \} \setminus \{a, b\}$ with $p_{k}=1$.

 We already distributed $2$ roots for $P$ and $Q$ each, so there is $1$ left to give. Our constraints give $p_{c}+p_{d}+p_{t}=1$, $p_{c},p_{d},p_{t}\in \{ 0,1 \}$ and $q_{t}\in \{0,1\}$, $q_{c}, q_{d}\in \{ 1,2 \}$. We can make some inferences:

 
            \begin{itemize}
        \item $p_{c}\in \{ 0,1 \}$, $q_c \in \{ 1,2 \}$, $m_c=p_{c}-q_{c} \in \{ -2,-1,0 \}$. Since $m_c$ is odd, $m_{c} = -1$. Similarly, $m_{d}=-1$.
        \item $p_{t}\in \{ 0,1 \}$, $q_{t}\in \{ 0,1 \}$, $m_{t}=p_{t}-q_{t}\in \{ -1,0,1 \}$. Since $m_{t}$ is even, $m_{t}=0$.
            \end{itemize}
            At this point, we have $(m_{a},m_{b},m_{c},m_{d},m_{t})=(1,1,-1,-1,0)$. But this is exactly the distribution of multiplicities of $\frac{(x-a)(x-b)}{(x-c)(x-d)}$, which is forced from the beginning. Thus, the remaining multiplicity of each of $P$ and $Q$ must cancel into a single bin to create net multiplicity $0=1-1$, yielding

 
 \[f(x)=\frac{(x-a)(x-b)}{(x-c)(x-d)}\cdot \frac{x-k}{x-k}\]where $k=c,d,t$, and this creates a hole if there isn't one already. However, the problem is that if there is already a hole (i.e. $k=c,d$), then this factor $\frac{x-k}{x-k}$ changes nothing as $f$ was already undefined at $x=k$, so while it gives a different pair of polynomials, it does not provide a new function $f$.

 Cases $k=c,d$ are symmetric:

 WLOG, let $k=c$, and choosing $t=1,2,3,4,5$ determines the remaining $1\le a<b<c<d \le 5$. Each $t$ creates a unique pair of polynomials, so there are $2\cdot \binom{5}{1}=10$ pairs of polynomials. However, as established, puncturing a hole again on $k=c,d$ does nothing:
 \[f(x)=\frac{(x-a)(x-b)}{(x-c)(x-d)}\cdot \frac{x-c}{x-c}=\frac{(x-a)(x-b)}{(x-c)(x-d)}\cdot \frac{x-d}{x-d}\]So each $t$ creates exactly one new function regardless of whether $k=c$ or $k=d$, i.e. there are only $5$ new functions.

 Case $k=t$: 

 Sketching a quick diagram for $(\text{sign pattern of }f)$, $t$ cannot puncture $[a,b]$ nor $(c,d)$, meaning it must lie outside them. All other cases cases do work, since you can puncture $+$ intervals without affecting $-$ intervals.
 \[\begin{cases} t,[a,b],(c,d) \\ [a,b],t,(c,d) \\ [a,b],(c,d),t \end{cases}\]Hence, there are $3$ valid choices of $t$.

 In total, there are $10+3=\boxed{ \textbf{(E)}~13 }$ ordered pairs of polynomials $(P,Q)$ but only $\boxed{ 8 }$ different functions created by $f=\frac{P}{Q}$.

 


\end{enumerate}